\documentclass[10pt,reqno]{amsart}
\usepackage[margin=1.75in,letterpaper]{geometry}   
\usepackage{setspace}
%\onehalfspacing
\usepackage{graphicx}
\usepackage{tikz}
\usepackage{tikz-cd}

\usetikzlibrary{graphs, positioning}
\usepackage{comment}
% Basic math packages
\usepackage{amsmath, amssymb, amsfonts, enumerate}
\usepackage{amsthm} % Theorem environments
\usepackage{mathtools} % Extended math tools
%\usepackage{enumitem} % Better control over lists
\usepackage[colorlinks=true, pdfstartview=FitV, linkcolor=blue, citecolor=blue, urlcolor=blue]{hyperref}
\usepackage[capitalise]{cleveref}
\crefformat{equation}{(#2#1#3)}
\crefrangeformat{equation}{(#3#1#4--#5#2#6)}
\crefformat{enumi}{(#2#1#3)}
\crefrangeformat{enumi}{(#3#1#4--#5#2#6)}


\makeatother

%%% THEOREM STYLES --------------------------------------------------

\newcounter{intro}
\newtheorem{introthm}[intro]{Theorem}
\renewcommand{\theintro}{\Alph{intro}}
\newtheorem{theorem}[subsection]{Theorem}
\newtheorem{proposition}[subsection]{Proposition}
\newtheorem{lemma}[subsection]{Lemma}
\newtheorem{corollary}[subsection]{Corollary}

\theoremstyle{definition}
\newtheorem{example}[subsection]{Example}
\newtheorem{setting}[subsection]{Setting}
\newtheorem{definition}[subsection]{Definition}
\newtheorem{notation}[subsection]{Notation}
\newtheorem{chunk}[subsection]{}
\newtheorem*{ack}{Acknowledgements}
\newtheorem{conjecture}[subsection]{Conjecture}
\newtheorem{obs}[subsection]{Observation}

\theoremstyle{remark}
\newtheorem{remark}[subsection]{Remark}
\newtheorem*{claim}{Claim}
\newtheorem{question}[subsection]{Puestion}
\newtheorem{construction}[subsection]{Construction}

\setcounter{secnumdepth}{2}
\numberwithin{equation}{section}



% Custom math shortcuts (optional)
\newcommand{\R}{\mathbb{R}}
\newcommand{\N}{\mathbb{N}}
\newcommand{\Z}{\mathbb{Z}}
\newcommand{\Q}{\mathbb{Q}}
\newcommand{\C}{\mathbb{C}}
\newcommand{\qwhere}{\quad \mbox{ where } \quad }
\newcommand{\qand}{\quad \mbox{ and } \quad }
\newcommand{\qfor}{\quad \mbox{ for } \quad }
\newcommand{\qforeach}{\quad \mbox{ for each } }
\newcommand{\qor}{\quad \mbox{ or } \quad }

\newcommand{\qif}{\quad \mbox{if} \quad}
\newcommand{\qwith}{\quad \mbox{with} \quad}
\newcommand{\qwhen}{\quad \mbox{when} \quad}
\newcommand{\qforsome}{\quad \mbox{for some} \quad}
\newcommand{\qforall}{\quad \mbox{for all} \quad}
\newcommand{\depth}{\mathrm{depth}}
\newcommand{\Tor}{\mathrm{Tor}}
\newcommand{\Ext}{\mathrm{Ext}}
\newcommand{\Soc}{\mathrm{Soc}}
\newcommand{\Exp}{\mathrm{Exp}}
\newcommand{\Log}{\mathrm{Log}}
\newcommand{\rate}{\mathrm{rate}}
\newcommand{\reg}{\mathrm{reg}}
\newcommand{\edim}{\mathrm{edim}}
\newcommand{\slant}{\mathrm{slant}}
\newcommand{\slope}{\mathrm{slope}}
\newcommand{\kk}{\mathsf k}
\newcommand{\Po}{\mathsf P}
\newcommand{\bo}{\mathsf b}
\newcommand{\Qo}{\mathsf Q}
\newcommand{\pd}{\mathrm{pd}}
\newcommand{\HH}{\mathrm{H}}
\newcommand{\Ker}{\mathrm{Ker}}
\newcommand{\m}{\mathfrak m}
\newcommand{\n}{\mathfrak n}
\newcommand{\st}{\colon\,}
\DeclareMathOperator*{\maximize}{maximize}
\DeclareMathOperator*{\subjectto}{\mbox{subject to}}

% Side comments
\newcommand{\nt}[2][$^\spadesuit$]{\hspace{0pt}#1\marginpar{\tt\raggedleft
    #1 #2}}
\newcommand{\adela}[2][$^\spadesuit$]{\nt[#1]{AV:#2}}
\newcommand{\dinesh}[2][$^\heartsuit$]{\nt[#1]{DL:#2}}
\numberwithin{equation}{section}


\title[Product of ideals and Golod threshold]{Golod thresholds of ideals  and products of ideals in dimension 3}


\author[D.~Limbu]{Dinesh Limbu}
\address{Department of Mathematics, University of South Carolina, 
Columbia, SC 29208}
\email{dlimbu@email.sc.edu}

\author[A.~Vraciu]{Adela Vraciu}
\address{Department of Mathematics, University of South Carolina, 
Columbia, SC 29208}
\email{vraciu@math.sc.edu}


\date{\today}

\begin{document}


\begin{abstract}
Let $I$ be a proper ideal in a standard graded polynomial ring $P$ with homogeneous maximal ideal $\m$. We show that there exists $k_0\in \mathbb{N}$ such that $\m^kI$ is Golod for all $k\geq k_0$. The smallest such $k_0$ is denoted $\mathrm{gt}_{\m}(I)$. We give an explicit formula for $\mathrm{gt}_{\m}(I)$ in the case when $I$ is a monomial complete intersection and bounds for $\mathrm{gt}_{\m}(I)$ for other classes of ideals. We show that under a certain degree condition the product of a Golod ideal with a maximal homogeneous ideal is Golod. In the last section, we show that in a regular local ring of dimension 3, the product of any two proper ideals is Golod; this solves a question posed by Dao-Stefani \cite[Question 4.5]{dao2022monomial} and VandeBogert \cite[Question 4.11]{vandebogert2022products}.
\end{abstract} 

\keywords{Golod ring, Poincar\'{e} series, homogeneous ideal, monomial complete intersection, Gorenstein ideal, trivial Massey operation}


\subjclass[2020]{13D02, 13D07, 13C13.} 

\maketitle

\section{Introduction}
Let $(P,\m,\kk)$ be a regular local ring of dimension $n$ over a field $\kk$ or a standard graded polynomial ring in $n$ variables, with homogeneous maximal ideal $\m$. Let $I$ be a proper ideal of $P$ and let $R=P/I$. The Poincar\'{e} series of $\kk $ over $R$ is $P_\kk^R(t)=\sum_{i\geq0}\beta_i^R(\kk)t^i$, where $\beta_{i}^R(\kk)=\dim_{\kk}[\Tor_i^R(\kk,\kk)]$, and the Poincar\'{e} series of $R$ over $P$ is $P_R^P(t)=\sum_{i\geq0}\beta_i^P(R)t^i$, where $\beta_{i}^P(R)=\dim_k[\Tor_i^P(R,\kk)]$. Serre established the following coefficient-wise inequality:
\[
P_\kk^R(t)
\leq
\frac{(1+t)^n}
{1-t\bigl(P_R^P(t)-1\bigr)}.
\]
When the above upper bound is attained, we say the ring $R$ is Golod and the ideal $I$ is Golod. The question of whether the product of two proper ideals is Golod has been studied in several papers: \cite{herzog1979golodideale}, \cite{herzog2013ordinary}, \cite{de2016products}, \cite{herzog2016homology}, \cite{maleki2019golod}, \cite{herzog2018koszul} and \cite{vandebogert2022products}. We briefly review some of the known results on this question. In the setting of a standard graded polynomial ring of characteristic $0$, Herzog and Huneke \cite{herzog2013ordinary} introduced strongly Golod ideals and showed that powers and symbolic powers of graded ideals are strongly Golod and the product of two strongly Golod ideals is Golod. Later, Herzog and Maleki \cite{herzog2018koszul} introduced $d$-Golod ideals in arbitrary characteristic and showed that the product of two $d$-Golod ideals is Golod. However, due to the work of De Stefani, we know that in more than three variables there are ideals whose product is not Golod. The following is the famous counterexample:\begin{example}\cite[Example 2.1]{de2016products}
    Let $P=\kk[x_1,x_2,x_3,x_4]$, $I=(x_1^2, x_2^2, x_3^2, x_4^2)$ and $\m=(x_1, x_2, x_3, x_4)$. Then $\m I$ is not Golod. 
\end{example} 
In the above example, we notice that $\m^kI$ is Golod for all $k\geq 2$. To capture this phenomenon, we define an invariant which we call the Golod threshold, as follows. 
\begin{definition}\label{golodthresholddef}
    Let $I$, $J$ be ideals in $P=\kk[x_1,\cdots,x_n]$.
    The {\em Golod threshold of $I$ with respect to $J$} is
    $$
    \mathrm{gt}_J(I):=\mathrm{min}\{k \in \mathbb N\, | \, J^{\ell } I \ \mathrm{is \ 
    Golod\ for \ all \ } \ell \ge k\},
    $$
    provided that such a $k$ exists.
\end{definition}
In Section 3 we show that the Golod threshold always exists: see \cref{existence of golod threshold}. We will then focus on the case $J=\m$ and establish upper bounds for all homogeneous ideals $I$ [\ref{another bound}, \ref{upper bound}] and a lower bound for homogeneous Gorenstein ideals $I$ [\ref{lower bound for gorenstein}], based on the graded shifts in the minimal free resolution of the ideal over the polynomial ring. Using these bounds, we compute the Golod threshold with respect to $\m$ for certain Gorenstein ideals; see Example \ref{gorenstein example}. At the end of this section, under a suitable condition on the degrees of the minimal generators of a homogeneous ideal $I$, we show that if $I$ is Golod then $\m I$ is Golod; see \cref{equigenratedgolodThreshold}. As a consequence, under suitable degree conditions, $\m^k I$ is Golod if and only if $k\geq \mathrm{gt}_{\m}(I)$.  

In Section 4, we develop a combinatorial criterion for the existence of a monomial basis for the Koszul homology; see \cref{shellmonbasis}. We apply this criterion to show that ideals of the form $\m^k (x_1^{a_1},\cdots,x_n^{a_n})$ admit monomial bases for their Koszul homology; see \cref{monomialbasisforCI}. 

In Section 5, we establish an exact formula for the Golod threshold with respect to $\m$ for the monomial complete intersection $(x_1^{a_1},\ldots,x_n^{a_n})$; see \cref{preciseCIbound}. 



In the last section, we focus on the case of dimension three, where several results are already known. Dao and De Stefani \cite{dao2022monomial} showed that in a standard graded three-variable polynomial ring, the product of any two proper monomial ideals is Golod. More recently, VandeBogert \cite{vandebogert2022products} showed that in a regular local ring of dimension three, the product of the maximal ideal with any proper ideal is Golod. The question of whether the product of two proper ideals in a regular local ring of dimension three is Golod had remained open. In \cref{maintheorem1} we prove that in a regular local ring of dimension three, the product of any two proper ideals is Golod. This resolves the questions of Dao-Stefani\cite[Question 4.5]{dao2022monomial} and VandeBogert\cite[Question 4.11]{vandebogert2022products} in the positive direction. 

% It all goes back to Herzog and Steurich \cite{herzog1979golodideale}{}, whose results imply that, in the regular local setting, the product of two transversal ideals $I$ and $J$ (i.e., $I\cap J=IJ$) is Golod. In \cite{huneke2007ideals}, Huneke showed that, over an unramified local ring, a Gorenstein ideal of height at least 2 can never be a product of two ideals. Since a non-hypersurface Gorenstein ideal is the opposite of a Golod ideal, Huneke's result, in some sense, suggested that products of ideals are Golod. In the setting of standard graded algebra with characteristic $0$, using the description of the Koszul cycle in terms of the partial derivatives, Herzog and Huneke \cite{herzog2013ordinary} show that if $\partial(I)^2\subseteq I$, then the ideal I is Golod; they named such ideals strongly Golod. Using this, they showed that the product of any two strongly Golod ideals is Golod. But later, Alessandro De Stefani gave a counterexample[\cite{de2016products}, Example 2.1] in 4 variables and showed that the product of two homogeneous ideals need not be Golod. More work has been done significantly in the $ 3$-variable setting. In \cite{dao2022monomial}, Dao and Stefani proved that the product of proper monomials in the $3$-variable standard graded algebra is always Golod, and recently in \cite{vandebogert2022products}, Vanderbogert proved that the product of the maximal ideal and any proper ideal in the regular local ring of dimension 3 is Golod. The following question is still open.
% \begin{question}
% Let $P=\kk[x_1,x_2,x_3]$ be a standard graded algebra. Is the product of any two proper ideals $ I$, $J$ in $R$ always Golod? 
% \end{question}
% Parallel to this, the Golodness of powers of ideals has also been studied. It is known that in the standard graded setting, $I^k$ is Golod in characteristic $0$ for all $k\geq 2$ \cite{herzog2013ordinary}, and for monomial ideals in any characteristic \cite{herzog2018koszul}; the non-monomial case in positive characteristic remains open.
% \begin{question}
%     Let $P$ be a standard graded algebra of embedding dimension at least 3, in positive characteristic $p$. If $I$ is a proper homogeneous non-monomial ideal. Is it true that $I^k$ is Golod for all $k \geq 2$?
% \end{question}

% Motivated by the above question, we started investigating the Golod property of ideal that look like the following: $J^kI$, where I and J are proper homogeneous ideals in a standard graded setting and $k\geq 1$. 
% \begin{example}[\cite{de2016products}, Example 2.1] In $P=k[x_1,\cdots,x_4]$, let $\m=(x_1,\cdots,x_4)$ and $I=(x_1^2,\cdots,x_4^2)$. Then the product ideal $\m I$ is not Golod. In the Koszul homology of $P/\m I$ one has
% \[
% [x_1x_2dx_1dx_2]\cdot
% [x_3x_4dx_3dx_4]
% =
% [x_1x_2x_3x_4dx_1dx_2dx_3dx_4]
% \neq0.
% \]    
% \end{example}
% % In the above example, one can check that multiplying another copy of $\m$ produces a Golod ideal. 
% We prove in \cref{existence of golod threshold} that for a fixed homogeneous ideal $I$, there exists an exponent $k_0$, which we will refer to as the {\em Golod threshold} of $I$, such that $m^kI$ is Golod for all $k\ge k_0$. In general, for any ideal $I$ and $J$, we can define the Golod threshold for $I$ over $J$ as the minimum exponent $k_0$ such that for all $k\geq k_0$, $J^kI$ is Golod. 
% We provide an upper bound for this invariant in terms of the degree of the generators and the regularity of $I$ (see \cref{another bound}). We also give a precise formula in \cref{preciseCIbound} for the Golod threshold of a monomial complete intersection ideal $I=(x_1^{a_1}, \cdots, x_n^{a_n})$.


% Along these lines, it would be useful to know that once $m^{k_0}I$ is Golod for one value of the exponent, then $m^{k}I$ continues to be Golod for $k\ge k_0$. However, this turns out not to be true. We provide a counterexample. Our counterexample also shows that the product of two Golod ideals need not be Golod.

% The main tool that we will use in calculating the Golod threshold of a monomial complete intersection ideal is the existence of a monomial cycle basis for the Koszul homology. Such a monomial basis is known to exist for monomial ideals with linear quotients and nice
% Koszul lifting property\cite{zoromski2024monomial} and this class includes stable monomial ideals, symmetric shifted ideals, and matroidal ideals. In this paper, we provide a new class of ideals for which the Koszul homology admits a monomial cycle basis, which includes ideals of the form $m^k(x_1^{a_1}, \ldots, x_n^{a_n})$.



\begin{comment}
This paper is inspired by the question of when a product of two proper homogeneous ideals is Golod. It is known that such products are not always Golod, and in fact there are many counterexamples where one of the ideal is the irrelevant ideal $m=(x_1, \ldots, x_n)$. However, we prove in \ref{theorem...} that for a fixed homogeneous ideal $I$, there exists an exponent $k_0$, which we will refer to as the {\em Golod threshold} of $I$, such that $m^kI$ is Golod for all $k\ge k_0$. In general, for any ideal $I$ and $J$, we can define the Golod threshold for $I$ over $J$ as the minimum exponent $k_0$ such that for all $k\geq k_0$, $J^kI$ is Golod. 
We provide an upper bound for this invariant in terms of the degree of the generators and the regularity of $I$ (see \ref{theorem...}). We also give a precise formula for the Golod threshold of a monomial complete intersection ideal $I=(x_1^{a_1}, \cdots, x_n^{a_n})$.


Along these lines, it would be useful to know that once $m^{k_0}I$ is Golod for one value of the exponent, then $m^{k}I$ continues to be Golod for $k\ge k_0$. However, this turns out not to be true. We provide a counterexample. Our counterexample also shows that the product of two Golod ideals need not be Golod.

The main tool that we will use in calculating the Golod threshold of a monomial complete intersection ideal is the existence of a monomial cycle basis for the Koszul homology. Such a monomial basis is known to exist for... 
We provide a new class of ideals for which the Koszul homology admits a monomial cycle basis, which includes ideals of the form $m^k(x_1^{a_1}, \ldots, x_n^{a_n})$.

\end{comment}


\section{Preliminaries}
Let $P=\kk[x_1,\ldots,x_n]$ be a standard multigraded polynomial ring over a field $\kk$ and $\m=(x_1,\ldots,x_n)$ be the maximal homogeneous ideal in $P$. Let $I$ be a proper ideal of $P$ and let $R=P/I$.
% The Koszul complex is one of the fundamental constructions in commutative algebra and has numerous applications throughout the subject. For more details, see \cite{avramov1998infinite}.
The Koszul complex on $x_1,\cdots,x_n$ over $R$ is \[
\mathbf{K}^{R}(\mathbf{x})=R\langle dx_1,\cdots,dx_n|\;\partial(dx_i)=x_i\rangle
\]
where $dx_i$ are exterior variables of degree 1. For $F=\{i_1<\cdots<i_r\}\subseteq [n]$, we will denote $dx_{i_1}\cdots dx_{i_r}$ by $dx_F$. For simplicity, we will use the notation $\mathbf {K}^{R}:=\mathbf{K}^{R}(\mathbf{x})$ and we will refer to this as the Koszul complex of $R$. For a regular local ring we will use the same notation as above, but instead of variables we will use the minimal generators of the maximal ideal to define the Koszul complex. For an element $z\in \mathbf{K}_i^R$, $|z|:=i$ denotes the homological degree of $z$.

\begin{chunk}
    If $I$ is a monomial ideal, the multigrading of $P$ naturally induces a multigrading on $R$, which in turn induces a multigrading on the Koszul complex of $R$. For $\alpha\in \mathbb{N}^n$, the $\alpha$-multigraded piece of $\mathbf{K}^{R}(\mathbf{x})$, would be\[
    [\mathbf{K}^R(\mathbf{x})]_{\alpha}=\bigoplus_{\substack{F\subseteq\mathrm{Supp}(x^{\alpha})\\ x^{\alpha-1_F}\notin I}}\kk x^{\alpha-1_F}dx_F
 \] 
 with the differential
 \[
\partial(x^{\alpha-1_F}dx_F)=\sum_{\substack{i\in F\\ x^{\alpha-1_{F\setminus \{i\} }}\notin I }} (-1)^{\mathrm{pos}_F(i)-1}x^{\alpha-1_{F\setminus \{i\} }}dx_{F\setminus \{i\} }
\]
where $\mathrm{pos}_F(i)$ is the position of $i$ in $F$ in the natural order and for $G\subseteq [n]$, $1_{G}=(g_1,\cdots,g_n)$ where $g_i=1$ if $i\in G$ otherwise $g_i=0$. 
\end{chunk}
\begin{chunk}
     Here we recall some terminology related to simplicial chain complexes. For more details, we refer to \cite[Chapter 1]{miller2005combinatorial}. A collection $\Delta\subseteq 2^{[n]}$ is called a simplical complex on vertex set $[n]$ if $F'\subseteq F$ and $F \in \Delta$ implies $F'\in \Delta$. The elements of $\Delta$ are called faces, and the maximal faces are called facets. The dimension of a face $F$ is $\dim F=|F|-1$. 
      Let \[
    \mathbb{F}_i=\{F\in \Delta:\;|F|=i+1\},\quad \mathbb{F}_{-1}=\{\emptyset\} \]
    and let $\kk^{\mathbb{F}_i}$ be the $\kk$-vector space with basis $\mathbb{F}_i$. The reduced chain complex of $\Delta$ over $\kk$ is \[
        \widetilde{\mathbf{C}}_{\cdot}(\Delta;\kk):\quad
0 \longrightarrow \mathbf{\kk}^{\mathbb{F}_{n-1}}
\xrightarrow{}
\cdots
\xrightarrow{}
\mathbf{\kk}^{\mathbb{F}_{0}}
\xrightarrow{}
\mathbf{\kk}^{\mathbb{F}_{-1}}
\longrightarrow 0.
\]
where for $F\in \mathbb{F}_i$ we define the differential by $
\partial_i(F)=\sum_{j\in F}(-1)^{\mathrm{pos}_{F}(j)-1}(F\setminus \{j\})
$. The $i^{th}$ reduced homology of $\Delta$ over $\kk$ is defined as $
\widetilde{H}_i(\Delta;\kk)=\ker(\partial_i)/\mathrm{im}(\partial_{i+1}).$ We will denote the reduced cochain complex of $\Delta$ over $\kk$ by $\widetilde{\mathbb C}^{\cdot}(\Delta;\kk)$, which is defind as $ \widetilde{\mathbb C}^{\cdot}(\Delta;\kk)
:=\mathrm{Hom}_{\kk}(\widetilde{{C}}_{\cdot}(\Delta;\kk),\kk)$ and its $i$th reduced cohomology of $\Delta$ over $\kk$ is defined as $\widetilde{H}^{i}(\Delta;\kk)=\ker(\partial^{i})/\mathrm{im}(\partial^{i-1})$. For a face $F\in \widetilde{\mathbf{C}}_{\cdot}(\Delta;\kk)$, we denote its dual in $ \widetilde{\mathbf{C}}^{\cdot}(\Delta;\kk)$ by $F^*$. 
     
\begin{comment}
\begin{definition}[Shellable Simplicial Complex]

A d-dimensional simplicial complex is called 'pure' if its maximal simplices all have dimension 'd'. Let $\Delta$ be a finite or countably infinite simplicial complex. An ordering $C_1,C_2,\ldots$ of the maximal simplices of $\Delta$ is a 'shelling' if, for all $k=2,3,\ldots$, the complex :\[
B_k:=\Big(\bigcup_{i=1}^{k-1}C_i\Big)\cap C_k
\]
is pure and of dimension one smaller than $\dim C_k$.
\end{definition}

\begin{definition}[Shifted Simplicial Complex]
A simplicial complex on $n$ vertices is shifted if there exists a labeling of the nodes by $1$ through $n$ such that for any face $\{v_1, v_2, \cdots, v_d\}$, replacing any $v_i$ by a node
with a smaller label results in a collection which is also a face. 
\end{definition}
\end{comment}
\end{chunk}

\begin{chunk} We will now recall an equivalent characterization of Golodness. For more details on Golod rings, one can check \cite{avramov1998infinite}.
 Let $I$ be a proper or proper homogeneous ideal in a regular local ring or a standard graded ring $P$ with maximal or maximal homogeneous ideal $\m$, and let $R=P/I$. Let $\mathfrak B$ be a $k$-basis for $H_{\geq 1}(\mathbf{K}^{R}(\mathbf{x}))$. A trivial Massey operation on the Koszul complex of $R$ is a function \[
     \mu:\bigsqcup _{q\geq 1} \mathfrak B^q\xrightarrow[]{}\mathbf{K}^{R}(\mathbf{x})  \ \ \text{such that}
     \]
\begin{itemize}
         \item[(i)] For $h\in \mathfrak B$, $\mu(h)$ is a cycle representing the class of $h$.
         \item[(ii)] For $q\geq 2$ and $h_1,\cdots,h_q\in \mathfrak B$, we can choose $\mu(h_1,\cdots,h_q)$ such that,
         \begin{equation}\label{massey}
         \partial \mu(h_1,\cdots,h_q)=\sum_{j=1}^{q-1}\overline{\mu(h_1,\cdots,h_{j})}\mu(h_{j+1},\cdots,h_q)
         \end{equation} 
         where $\overline{a}=(-1)^{|a|+1}a$.
         Note that $|\mu(h_1, \ldots, h_q)|=\sum_{j=1}^q |h_i|+q-1$.
         \item[(iii)] $\mathrm{im}(\mu)\subseteq \m \mathbf{K}^R(\mathbf{x})$
     \end{itemize}
 The ideal $I$ is Golod if and only if $\mathbf{K}^{P/I}(\mathbf{x})$ admits a trivial Massey operation.
 \begin{remark}\label{notGolodcondition}
 If $\mathbf{K}^R({\bf x})$ admits a trivial Massey operation, then the product of any two elements in $H_{\ge 1}({\bf K}^R({\bf x}))$ is zero.
 \end{remark}  
 \begin{remark}\label{zeroproduct} Suppose that we can choose a basis $\mathfrak B$ for $H_{\geq 1}(\mathbf{K}^{R}(\mathbf{x}))$, such that for every $h\in \mathfrak B$, there is a cycle $z_h$, with $h=[z_h]$ and \[z_hz_{h'}=0\quad\text{for all }h,h'\in \mathfrak B.
 \]Then we may define $\mu(h)=z_h$ and $\mu(h_1,\cdots,h_p)=0$ for $p\geq 2$. One can check that this is indeed a trivial Massey operation. The following is a slightly more general version of this observation.
 \end{remark}
\end{chunk}


\begin{lemma}\label{trivialmassey1}
Let $p:=\mathrm{pd}_P(P/I)$. Assume that $\mathfrak B$ is a basis for $H_{\ge 1}(\mathbf{K}^{R}(\mathbf{x}))$, and for every $h\in \mathfrak B$ there is a cycle $z_h$ with $[z_h]=h$ such that $z_hz_{h'}=0$ for all $h, h'\in \mathfrak B$ with $|h|+|h'|\le p$.
 Then $I$ is Golod.
\end{lemma}
\begin{proof}
 We can define $\mu(h)=z_h$ and $\mu(h_1, \ldots, h_q)=0$ if $q\ge 2$ and $|h_1|+\cdots + |h_q|\le p$. For $|h_1|+\cdots + |h_q|>p$, $\mu(h_1, \ldots, h_q)$ can be defined recursively by noting that once $\displaystyle \mu: \bigsqcup_{1\le q'\le q-1}\mathfrak B^{q'}\to \mathbf{K}^{R}(\mathbf{x})$ has been defined, the expression on the right hand side of equation (\ref{massey}) is a cycle in homological degree $>p$; since $H_{>p}(\mathbf{K}^{R}(\mathbf{x}))=0$, we can choose $\mu(h_1, \ldots, h_q)$ such that (\ref{massey}) is satisfied.
\end{proof}

\begin{corollary}
If $I\subseteq \m^2$ is an ideal with $\mathrm{pd}_P(P/\m I)\le 3$, then $\mathfrak m I$ is Golod.  
\end{corollary}
\begin{proof}
Let $R=P/\mathfrak m I$. Due to the isomorphism $H_1(\mathbf{K}^{R}(\mathbf{x}))\cong \mathrm{Tor}_1^P(P/\m I, P/\mathfrak m)$, we can choose a basis for $H_1(\mathbf{K}^{R}(\mathbf{x}))$
consisting of elements of the form $[gdx_i]$ with $g\in I$ (corresponding to the generators $x_ig$ of $\mathfrak m I$).
 Since every cycle $z_h$ has coefficients in $\m$, it follows that $z_hz_{h'}=0$ for any $h, h'$ with $|h|+|h'|\le 3$.
 
 % Let $z_h=\sum_{|F|=i} g_Fdx_F$, then $\partial(z_h)=0$ implies $\sum_{\ell \notin G} (-1)^{\mathrm{pos}_{G\cup\{\ell\}}(\ell) }{x_{\ell}  g_{G\cup \{\ell\}}}=0$ in $P/I$, now if one of $g_{G\cup \{\ell\}}$ is unit then $\sum_{\ell \notin G} (-1)^{\mathrm{pos}_{G\cup\{\ell\}}(\ell) }{x_{\ell}  g_{G\cup \{\ell\}}}  $ is not zero in $P/\m^2$ which is a contraction.
 \end{proof}



 
% \begin{corollary}
%     In $P=k[x_1,x_2,x_3,x_4]$, let $I$ be a proper ideal, then $\m I$ is Golod if $\m$ is not an Associated prime of $P/I$.
% \end{corollary}

% In $P=k[x_1,x_2,x_3,x_4]$, if $\depth (P/I)=0$ and $\text{height}(I)\leq 3$, then $\m I$ is Golod?\\
% No, we have $(x_1^3,x_2^3,x_3^3,x_1x_4^2)$.\\

% In $P=k[x_1,x_2,x_3,x_4]$, if $\depth (P/I)=0$ and $\text{height}(I)\leq 2$, then $\m I$ is Golod?\\
% No, we have $(x_1^3,x_2^3,x_2x_3^2,x_1x_4^2)$.


\section{Existence of Golod Threshold with Some Upper and Lower Bounds.}
% In this section, we will introduce the notion of Golod threshold and show that it always exists; establish some general upper bound for all ideals and a lower bound for Gorenstein ideals. At the end, using these bounds, we will compute the Golod threshold for a class of Gorenstein ideals. 
Let $I$ and $J$ be two proper ideals in the standard graded polynomial ring $P=\kk[x_1,\cdots,x_n]$ with maximal homogeneous ideal $\m$ (or a regular local ring $(P,\m,\kk)$). In this section, we show that the Golod threshold of $I$ with respect to $J$ exists, and then we will focus on the case $J=\m$.
 The result below follows the steps in [Lemma 3.1, \cite{maleki2019golod}], with a slight modification.
\begin{lemma}\label{artinrees1}
    Let $M$ be a finitely generated $P$-module, and $N$ be a submodule of $M$. For any proper ideals $I$ and $J$ in $P$, there exists $d_0\in \mathbb{N}$ such that for all $d\geq d_0$ we have \[
    (J^dIM)\cap N=J^{d-d_0}(J^{d_0}IM\cap N).
    \]
\end{lemma}
\begin{proof}
    The Artin-Reese Lemma for $N\cap IM$ as a submodule of $IM$ implies that there exists $d_0$ such that for all $d\geq d_0$,\begin{equation}\label{artinrees2}
    J^dIM\cap (N\cap IM)=J^{d-d_0}(J^{d_0}IM\cap (N\cap IM))
    \end{equation}
    Observe that left hand side of \cref{artinrees2} is $J^dIM\cap N$ and the right hand side is $J^{d-d_0}(J^{d_0}IM\cap N)$.
\end{proof}

\begin{proposition}\label{torgolod}
     Let $I$ and $J$ be ideals in $P$. There exists an integer $d_0$ such that for all $d>d_0$, the natural map \begin{equation}\label{eqtorgolod}
    \Tor_i^P(P/J^dI, \kk)\xrightarrow[]{}\Tor_i^P(P/J^{d-1}I, \kk)
    \end{equation}
    induced by the surjection $P/J^dI\xrightarrow[]{}P/J ^{d-1}I$ is zero for all $i>0$.
\end{proposition}
\begin{proof}
     Set $\mathbf K^P=\mathbf{K}^P(\mathbf{x})$. Let $Z$ and $B$ be the cycles and boundaries in $\mathbf K^P$, respectively. Consider
    \[
    H_{*}(J^dI\otimes_P \mathbf K^P)\xrightarrow[]{\;\;\;\;}H_{*}(J^{d-1}I\otimes_P \mathbf K^P)
    \] 
 The map \cref{eqtorgolod} is zero if and only if \[\ker(J^dI\otimes_P \mathbf K^P_i\xrightarrow[]{}J^dI\otimes_P \mathbf K^P_{i-1})\subseteq \mathrm{im}(J^{d-1}I\otimes_P \mathbf K^P_{i+1}\xrightarrow[]{}J^{d-1}I\otimes_P \mathbf K^P_{i}),\]
 i.e., $J^dI\mathbf K^P_i\cap Z_i\subseteq J^{d-1}IZ_i$. Now, \cref{artinrees1} implies that there exists $d_0$ such that for all $d\geq d_0$, \begin{equation}\label{eq1}
 J^dI\mathbf K^P_i\cap Z_i=J^{d-d_0}(J^{d_0}I\mathbf K^P_i\cap Z_i)   
\end{equation}
for all $i\geq 0$. Since $H_i(J^{d_0}I\mathbf K^P)$ is a $\kk$-vector space we have \begin{equation}\label{eq2}
J(J^{d_0}I\mathbf K^P_i\cap Z_i)\subseteq B_{i}(J^{d_0}I\mathbf K^P)=J^{d_0}IZ_i
\end{equation}
which implies \[
 J^dI\mathbf K^P_i\cap Z_i=_{\eqref{eq1}}J^{d-d_0-1}J(J^{d_0}I\mathbf K^P_i\cap Z_i)\subseteq_{\eqref{eq2}} J^{d-d_0-1}( J^{d_0}IZ_i)=J^{d-1}IZ_i
\] 
\end{proof}

\begin{theorem}\label{existence of golod threshold}
    $\mathrm{gt}_J(I)$ exists for all proper ideals $I$ and $J$ of $P$.
\end{theorem}
\begin{proof}
    Let $L=J^{k-1}I$, where $k\geq 2$. Since $L^2\subseteq J^kI\subseteq L$, from Lemma 2.1 in \cite{maleki2019golod} and \cref{torgolod}, there exists $k_0$ such that for all $k\geq k_0$, $J^kI$ is Golod. 
\end{proof}
 We will establish upper bounds for $\mathrm{gt}_{\m}(I)$ in terms of the regularity of $I$; we recall the definiton below: 
 \begin{definition}
     Let $I\subseteq P=\kk[x_1,\ldots,x_n]$ be a proper homogeneous ideal. The regularity of $I$ denoted by $\reg(I)$ is \[
    \reg(I):= \max_{1\leq i\leq n}\{j-i\;|\; \dim_{k}[\Tor^{P}_i(P/I,k)]_j\neq 0\}.\]
 \end{definition}
The following result, stated as Theorem 6.16 (b) in \cite{mccullough2015infinite}, was privately communicated by Huneke to the authors of \cite{mccullough2015infinite}. We could not find a published proof in the literature; we include one here for completeness. 

\begin{proposition}\label{huneke}
    Let $P=\kk[x_1,\ldots,x_n]$ be a standard graded algebra with homogeneous maximal ideal $\m$. Let $I$ be a proper homogeneous ideal of $P$, and let $d$ be the minimal degree among the generators of $I$. If $ \reg(I)\leq 2d-3$, then $I$ is Golod. 
\end{proposition}
\begin{proof}
    Let $R=P/I$ and $\overline{\m}$ be the homogeneous maximal ideal in $R$. The assumption implies that $I\subseteq \m^2$ (since $d\ge 2$) and for $i\geq 1$, we have: \[
    \dim_{\kk}[\Tor_{i}^P(R,\kk)]_j=0 \quad\text{if}\quad j-i<d-1 \quad\text{or}\quad j-i\geq 2(d-1).
    \] 
    The conclusion follows from \cite[Lemma 9.4]{diethorn2026homological} and \cite[Proposition 9.3]{diethorn2026homological}.
\end{proof}


% We know that $\reg(\m^kI)\leq k+\reg(I)$, so if $\reg(I)\leq 2d+k-3$, then $\m^kI$ is Golod. In other words, for all $k\geq \reg(I)-2d+3$, $\m^kI$ is Golod. Therefore, an upper bound for $\mathrm{gt}_\m(I)$ is $\reg(I)-2d+3$. But we can do better; below is a general upper bound for $\mathrm{gt}_m(I)$.
% \begin{corollary}\label{upperbound}
%     If $I$ is a homogeneous ideal, we have $\mathrm{gt}_{\m}(I)\le \reg(I)-2d+3$, where $d$ is the minimal degree among the generators of $I$.
% \end{corollary}
\begin{corollary}\label{another bound}
Let $I$ be a homogeneous ideal and let $d$ and $D$ denote the smallest and respectively the largest degree of a minimal generator of $I$. Then $$\mathrm{gt}_{\m}(I) \le \mathrm{max}\left\{\dfrac{\mathrm{reg}(I)-2d+3}{2}, D-2d+2\right\}.$$
\end{corollary}
\begin{proof}
Let $k\ge \mathrm{max}\{\frac{\mathrm{reg}(I)-2d+3}{2}, D-2d+2\}$. 
Using \cite[Corollary 7.5]{eisenbud2006regularity} and $\reg(\m^k)=k-1$, we have $$\mathrm{reg}(\m^kI)\le \mathrm{max}\{\mathrm{reg}(I), k+D-1\}, 
$$
and the smallest degree of a generator of $\m^kI$ is $k+d$. The assumption implies 
$\mathrm{reg}(\m^kI)\le 2(k+d)-3$, and the conclusion follows from Proposition (\ref{huneke}).
\end{proof}

The following result will allow us to obtain another upper bound for $\mathrm{gt}_{\m}(I)$:

\begin{proposition}\label{tor zero}
Let $I$ be a homogeneous ideal and let $p:=\mathrm{pd}_P(P/I)$. If $k$ is such that 
the maps $\mathrm{Tor}_i^P(P/\m^kI, \kk)\to \mathrm{Tor}_i^P(P/I, \kk)$ induced by the canonical projection $P/\m^kI\to P/I$ are zero for all $1\le i \le \lfloor p/2\rfloor$, then $\m^kI$ is Golod. 
\end{proposition}
\begin{proof}
Consider the short exact sequence 
\begin{equation}\label{ses}
0\to \frac{I}{\m^k I} \to \frac{P}{\m^kI}\to \frac{P}{I}\to 0.
\end{equation}
 Tensoring (\ref{ses}) with ${\bf K}^P({\bf x})$ yields short exact sequences
$$
H_i({\bf K}^P({\bf x})\otimes I/\m^k I) \to H_i({\bf K}^{P/\m^kI}({\bf x}))\to H_i({\bf K}^{P/I}({\bf x})).
$$
Due to the isomorphism between Tor and Koszul homology, the rightmost map is zero for $1\le i \le \lfloor p/2 \rfloor$ by assumption. This implies that $H_i({\bf K}^{P/\m^kI}({\bf x}))$ is equal to the image of the map 
$H_i({\bf K}^P({\bf x})\otimes I/\m^kI)\to H_i({\bf K}^{P/\m^kI}({\bf x}))$. 

    % ({\em note: commenting out the thing about coefficients in $\m^{k-1}I$})
    % We claim that the assumption that the ideal is equigenerated implies that all such cycles must in fact have coefficients in $\m^{k-1}I$\dinesh{This also means for $i\leq \lfloor p/2 \rfloor $ we have monomial basis for $H_i({\bf K}^P\otimes P/\m^kI) $}. Let $d$ denote the degree of the generators of $I$. An element in $H_i({\bf K}^P\otimes P/m^kI)$ must have degree at least $d+k+i-1$. If $[z]=\sum_{|F|=i} g_F dx_F$ is homogeneous of degree $\ge d+k+i-1$, then $\mathrm{deg}(g_F)\ge d+k-1$ and, given $g_F\in I$, this implies $g_F\in \m^{k-1}I$.

It follows that for all $1\le i\le \lfloor p/2 \rfloor$, every element of $H_i({\bf K}^{P/\m^kI}({\bf x}))$ can be represented by a cycle that has coefficients in $I/\m^kI$.
Also note that for all $i\ge 1$, every nonzero element in $H_i({\bf K}^{P/\m^kI}({\bf x}))$ must have degree at least $k+i$, and therefore it is represented by a cycle with coefficients in $\m^k/\m^kI$. 
    
    Consider a basis $\{[z_{\lambda}]\}$ for $H_{\ge 1}({\bf K}^P({\bf x})\otimes P/\m^kI)$ with $[z_{\lambda}]\in H_{i_{\lambda}}$. 
    Let $\lambda_1, \lambda_2$ be such that $i_{\lambda_1}+i_{\lambda_2}\le p$. 
 We may assume without loss of generality that $i_{\lambda_1}\le \lfloor \frac{p}{2}\rfloor.$
 Since $z_{\lambda_1}$ has coefficients in $I/\m^kI$ and $z_{\lambda_2}$ has coefficients in $\m^k/\m^kI$, it follows that $z_{\lambda_1}z_{\lambda_2}=0$ and the conclusion follows from
\cref{trivialmassey1}.
\end{proof}
\begin{corollary}\label{upper bound}
   Let $I$ be a homogeneous ideal with $p:=\mathrm{pd}_P(P/I)$, and $d$ be the smallest degree of a generator of $I$.
For each $i\le p$, let $$T(i):=\mathrm{max}\{j: [\Tor_{i}^P(P/I,\kk)]_j\ne 0\}.$$
   Then
$$
\mathrm{gt}_{\m}(I)\le \mathrm{max}\{T(i)-i-d+2: i \le \lfloor p/2 \rfloor \}
$$ 
\end{corollary}
\begin{proof}
Let $k\ge \mathrm{max}\{T(i)-i-d+2: i \le \lfloor p/2 \rfloor \}$.
  We have $[\mathrm{Tor}_i(P/I, \kk)]_j=0$ for $j>T(i)$, and 
  $[\mathrm{Tor}_i(P/\m^kI, \kk)]_j=0$ for $j<d+k+i-1.$
  The assumption implies that for all $i\le \lfloor \frac{p}{2}\rfloor$, we have $d+k+i-1>T(i)$, and therefore the map
  $$
[\mathrm{Tor}_i(P/I, \kk)]_j\to [\mathrm{Tor}_i(P/\m^kI, \kk)]_j
  $$
  is zero in each degree $j$. \cref{tor zero} shows that $\m^kI$ is Golod and the conclusion follows.
\end{proof}

\begin{example}
Let $I=(x_1^d,\ldots,x_c^d)$ be a monomial complete intersection in $P=\kk[x_{1},\cdots,x_{n}]$. When $c$ is even, both \cref{another bound} and \cref{upper bound} give the same upper bound for $\mathrm{gt}_m(I)$, namely $(\frac{c}{2}-1)(d-1)+1$; when $c$ is odd, \cref{upper bound} gives a better upper bound, namely $(\lfloor \frac{c}{2}\rfloor -1)(d-1)+1$. It turns out that the value of $\mathrm{gt}_{\m}(I)$ is in fact equal to this upper bound. We will calculate the value of $\mathrm{gt}_{\m}(I)$ for a general monomial complete intersection $I=(x_1^{a_1}, \ldots, x_n^{a_n})$ in \cref{preciseCIbound}.

% We have $\mathrm{pd}_P(P/I)=c$,  $T(i)=id$ for all $1\le i \le c$ and $\reg(I)=c(d-1) $. Suppose $d\geq 2$, then for $i\leq c$, we have \[
% \mathrm{max}\{T(i)-i-d+2: i \le \lfloor c/2 \rfloor \}=\left({\left\lfloor \frac{c}{2}\right\rfloor}-1\right)(d-1)+1
% \]
% and 
% \[
% \mathrm{max}\left\{\dfrac{\mathrm{reg}(I)-2d+3}{2}, D-2d+2\right\}=\left(\frac{c}{2}-1\right)(d-1)+1
% \]
% Therefore, when $c$ is even both \cref{another bound} and \cref{upper bound} give the same upper bound for $\mathrm{gt}_m(I)$, and when $c$ is odd \cref{upper bound} gives a better upper bound.

\end{example}

If the ideal $I$ is Gorenstein, the following provides a lower bound for $\mathrm{gt}_{\m}(I)$:

\begin{theorem}\label{lower bound for gorenstein}
Let $I\subseteq P$ be a homogeneous Gorenstein ideal of codimension $c$ and regularity $\mathrm{reg}(I)$. Let $b_{ij}$ denote the dimension of $[\mathrm{Tor}_i(P/I, \kk)]_j$. For $0\le i\le c$, let $t(i)$ and $T(i)$ denote the smallest and respectively largest values of $j$ such 
$b_{ij}\ne 0$ and let $D$ be the largest degree of a generator of $I$.

Let $(i_1, j_1)$ be such that $1\le i_1\le c-1$ and $b_{i_1 j_1}\ne 0$, and let $i_2=c-i_1$, $j_2=c+\mathrm{reg}(I)-j_1$.

Then 
$$
\mathrm{gt}_{\m}(I)> \mathrm{min}\{j_1-i_1-D+1, j_2-i_2-D+1\}
$$\end{theorem}
Before proving this theorem, we will need a Lemma:
\begin{lemma}\label{mod k}
Let $R=P/I$ and $R_k=P/\m^kI$ where $I$ is a homogeneous ideal of the polynomial ring $P$, and let $D$ denote the largest degree of a generator of $I$. Let $\pi_k: {\bf K}^{R_k}\to {\bf K}^R$ denote the map of complexes induced by  canonical projection $\pi: R\to R_k$.
For every $[z]\in H_i({\bf K}^{R})$ with
$\mathrm{deg}([z])\ge k+D+i-1$, there exists $[\tilde{z}]\in H_i({\bf K}^{R_k})$ such that $[z]=\pi_k(\tilde{[z]})$. 
\end{lemma}
\begin{proof}
Write $$[z]=[\sum_{|F|=i} \overline{g}_Fdx_F]$$
where the summation is over subsets $F\subseteq [n]$ and $g_F$ are elements of $P$ with $\overline{g}_F$ being their image in the quotient ring $R$. Let $\tilde{g}_F\in R_k$ be the image of $ g_F$ in the quotient ring $R_k$.
The assumption that $[z]\in H_i({\bf K}^R)$ means that $\sum_{|F|=i} \overline{g}_Fdx_F$ is a cycle in the Koszul complex, i.e. for each $G\subseteq [n]$ with $|G|=i-1$ we have 
$$
\sum_{\ell \notin G} (-1)^{\mathrm{pos}_{G\cup\{\ell\}}(\ell) }\overline{x_{\ell } g_{G\cup \{\ell\}}}=0,$$ 
and therefore 
$$
\sum_{\ell \notin G} (-1)^{\mathrm{pos}_{G\cup\{\ell\}}(\ell) }x_{\ell } g_{G\cup \{\ell\}}\in I.$$
The degree assumption implies that 
$$
\sum_{\ell \notin G} (-1)^{\mathrm{pos}_{G\cup\{\ell\}}(\ell)} x_{\ell } g_{G\cup \{\ell\}}\in \m^k I,$$
and therefore 
$\tilde{z}:=\sum_{|F|=i} \tilde{g}_F dx_F$ is a cycle in ${\bf K}^{R_k}$.
\end{proof}

We are now ready to prove \cref{lower bound for gorenstein}:
\begin{proof}
Let
\begin{equation}\label{lower bound}
k\le \mathrm{min}\{j_1-i_1-D+1, j_2-i_2-D+1\}
    \end{equation}
% $0\le i \le c$ 
% such that 
% \begin{equation}\label{lower bound2}
% k\le \mathrm{min}\{t(i)-i-D+1, \mathrm{reg}(I)-T(i)+i-D+1\}
% \end{equation}
We will prove that $\m^k I$ is not Golod.




Due to the isomorphism $\mathrm{Tor}_i^P(P/I, \kk)\cong H_i({\bf K}^R)$, $b_{i_1 j_1}\ne 0$ implies that there is a nonzero $[z_1]\in H_{i_1}({\bf K}^{R})$ with $\mathrm{deg}([z_1])=j_1$, and Poincar\'{e} duality implies that there exists a complementary $[z_2]\in H_{i_2}({\bf K}^R)$ with $\mathrm{deg}([z_2])=j_2$ such that $[z_1][z_2]\ne 0$ in $H_c({\bf K}^R).$
\cref{mod k}
yields
$[\tilde{z_1}]\in H_{i_1}({\bf K}^{R_k})$ with $\pi_k([\tilde{z_1}])=[z_1]$ and 
$[\tilde{z_2}]\in H_{i_2}({\bf K}^{R_k})$ with $\pi_k([\tilde{z_2}])=[z_2]$. Since $[z_1][z_2]\ne 0$, it follows that $\tilde{[z_1]}\tilde{[z_2]}\ne 0$, and therefore $\m^k I$ is not Golod.
\end{proof}

% {\em commenting out the way proof was originally written }

% Since $j\ge t(i)$, the assumption (\ref {lower bound}) implies $\mathrm{deg}([z])\ge k+D+i-1$, and Observation (\ref{mod k}) implies that there exists $[\tilde{z_1}]\in H_i(\tilde{\bf K}_k)$ with $\pi_k([\tilde{z_1}])=[z_1]$.   

% Poincar\'{e} duality implies that there exists $[z_2]\in H_{c-i}({\bf K})$ such that $[z_1][z_2]\ne 0$, and $\mathrm{deg}([z_2])=c+\mathrm{reg}(I)-j$. Since $j\le T(i)$, the assumption (\ref{lower bound}) implies that $\mathrm{deg}([z_2])\ge k+D+(c-i)-1$, and Observation (\ref{mod k}) yields a $\tilde{[z_2]}\in H_{c-i}(\tilde{\bf K}_k)$ such that $[z]_2=\pi_k(\tilde{[z_2]})$. Since $[z_1][z_2]\ne 0$, it follows that $\tilde{[z_1]}\tilde{[z_2]}\ne 0$, and therefore $\m^k I$ is not Golod.
\begin{example}\label{cyclicGorensteinexample}
For $n\geq 5$ and $a\geq 1$, consider the ideal 
$$
I_{a,n}=(x_i^ax_j^a:\; \ i, j \in [n] \ \text{and} \ |i-j|\neq 0,1,n-1)\subset P=\kk [x_1, \ldots, x_n]
$$
Then 
$$
\mathrm{gt}_{\m}(I_{a, n})=\left\lfloor \frac{n-4}{2}\right\rfloor  (a-1)+1
$$
\end{example}
\begin{proof} We will describe the pairs $(i, j)$ for which $b_{ij}\ne 0$, where $b_{ij}$ are graded betti numbers of $P/I_{a, n}$, and then apply \cref{upper bound} and \cref{lower bound for gorenstein}.
We start by observing that $P/I_{1, n}$ is the Stanley-Reisner ring of the cycle $C_n$ with vertices $\{1, \ldots, n\}$.
For a squarefree monomial $\displaystyle M=\prod_{i\in W} x_i$, Hochster's formula \cite[Theorem 5.1]{hochster1977cohen} for the multigraded Betti numbers gives
$$
b_{i, M}(P/I_{1, n})=\mathrm{dim}_{\kk} (\tilde{H}_{|W|-i-1}(\Delta_W, \kk))
$$
where $\Delta_W$ is the simplicial complex induced by $C_n$ on $W$.

For $|W|<n$, $\Delta_W$ is a disjoint union of paths and therefore 
$\tilde{H}_r(\Delta_W)=0$ for all $r>0$, and $\tilde{H}_0(\Delta_W)=c(W)-1$, where $c(W)$ denotes the number of connected components of $\Delta_W$. In particular, 
$
b_{i, M}(P/I_{1, n})\ne 0$ if and only if $i=\mathrm{deg}(M)-1$ and $\Delta_W$ is disconnected, i.e. $M$ is not a product of $i+1$ consecutive variables (where $x_n, x_1$ are considered consecutive).

For $W=[n]$ (i.e. $M=\prod_{i=1}^n x_i$), we have $\Delta_W=C_n$, and $\tilde{H}_r(C_n)\ne 0\Leftrightarrow r=1$; more precisely, $\tilde{H}_1(C_n)\cong \kk$. It follows that 
$b_{i, M}\ne 0\Leftrightarrow i=n-2$. 
We see that $\mathrm{pd}_P(P/I_{1, n})=n-2$ and
$b_{n-2, M}\ne 0\Leftrightarrow M=\prod_{i=1}^n x_i$, in which case $b_{n-2, M}=1$.

In view of Theorem 3.3 in \cite{gasharov1999lcm}, the multigraded Betti numbers in the resolution of $I_{a, n}$ are obtained from those in the resolution $I_{1, n}$ by replacing each monomial $M=\prod_{i\in W} x_i$ by $M_a:=\prod_{i\in W} x_i^a$. In terms of the standard grading, we have $b_{i, j}\ne 0 \Leftrightarrow j=(i+1)a$ for $i\le n-3$, and $b_{n-2, j}\ne 0\Leftrightarrow j=na$.

We apply \cref{upper bound} to obtain an upper bound for $\mathrm{gt}_{\m}(I_{a, n})$:
\begin{equation}\label{gorenstein example}
    \begin{split}
\mathrm{gt}_{\m}(I_{a, n})& \le \mathrm{max}\{ (i+1)a-i-2a+2  \, : \,   i\le \left\lfloor \frac{n-2}{2}\right\rfloor \}\\
&=\left\lfloor \frac{n-4}{2}\right\rfloor (a-1)+1\\
\end{split}
\end{equation}
Consider the ring homomorphism $\psi:P\xrightarrow[]{}P$, where $\psi(x_i)=x_i^a$. Here $\psi$ is faithfully flat, as $P$ is a free $P$-module via $\psi$. Changing the base from $P$ to $P/I_{1,n}$, we have the flat map $ \psi\otimes P/I_{1,n} :P/I_{1,n}\xrightarrow[]{}P/I_{a,n}.$ Now, since $P/I_{1,n}$ and $P/(x_1^a,\ldots,x_n^{a})$ are Gorenstein, after localizing at $\m$ and using \cite[Corollary 3.3.15]{bruns1998cohen}, we have $P/I_{a,n}$ is Gorenstein.

 We can now use \cref{lower bound for gorenstein} to see that have equality in equation \cref{gorenstein example}.
Indeed, let $(i_1, j_1)$, $(i_2, j_2)$ be as in the assumption of \cref{lower bound for gorenstein}, so that $i_1+i_2=n-2$, $j_1=(i_1+1)a$, $j_2=(i_2+1)a$. Without loss of generality, assume $i_1\le i_2$. We have
$$
\mathrm{gt}(I_{a, n})>\mathrm{min}\{(i_1-1)(a-1), (i_2-1)(a-1)\}=(i_1-1)(a-1)
$$
Since this holds for any choice of $i_1, i_2$ subject to $i_1+i_2=n-2$, we can take $i_1=\lfloor \frac{n-2}{2}\rfloor$, $i_2=\lceil \frac{n-2}{2}\rceil$ and the conclusion follows.
\end{proof}

The proof of \cref{cyclicGorensteinexample} shows that for the Gorenstein ideal $I_{a,n}$, if $m^kI_{a,n}$ is Golod then $m^{k'}I_{a,n}$ is also Golod for all $k'>k$. The following example shows that this behavior does not always hold.

\begin{example}
    Let $I=(x_1^3, x_2^3, x_3^3,x_4^3, x_1x_2^2x_3^2x_4^2)\subseteq \kk[x_1, \ldots, x_4]$. For this ideal, $\m I$ is Golod (one can check that the criterion \cite[Theorem 1.5]{zoromski2024monomial} is satisfied for $\m I$) but $\m^2I$ is not.  Indeed,    $[x_1^2x_2^2dx_{12}]$ and $[x_3^2x_4^2dx_{34}]$ are elements in $H_{\ge 1}(K^{P/\m^2I}(\mathbf x))$
    and their product is not zero. 
\end{example}

The above example also resolves Question 5.5 \cite{de2016products} in the negative, as $\m I $ and $\m$ are Golod, but $\m\cdot \m I$ is not Golod. Now, we note that even though there are examples where the Golod property of $\m^k I$ oscillates, there are an abundance of ideals for which the Golod threshold is sharp, i.e., if $\m^kI$ is Golod then $\m^{k'}I$ is Golod for all $k'\geq k$. To prove this, we will show that we can pull back the trivial Massey operation under a suitable degree condition. For a similar construction under different hypotheses, one can look at \cite{vandebogert2022detecting}.
 
 \begin{theorem}\label{equigenratedgolodThreshold}
     Let $I \subseteq \m^2$ be a homogeneous ideal in a standard graded polynomial ring $P$. Let $d$ and $D$ denote the smallest and largest degrees, respectively, of a minimal generator of $I$. If $D\leq 2(d-1)$ and $I$ is Golod, then $\m I$ is also Golod. 
 \end{theorem}
 
 \begin{proof}
 Consider the canonical short exact sequence
     \[
     0\xrightarrow[]{}\mathbf{K}^{I/\m I}(\textbf{x})\xrightarrow[]{} \mathbf{K}^{P/\m I}(\textbf{x})\xrightarrow[]{\varphi}\mathbf{K}^{P/I}(\textbf{x})\xrightarrow[]{}0,
     \]
     and the exact sequence of homology \[
      H(\mathbf{K}^{I/\m I}(\textbf{x}))\xrightarrow[]{} H(\mathbf{K}^{P/\m I}(\textbf{x}))\xrightarrow[]{\varphi_*}H(\mathbf{K}^{P/I}(\textbf{x}))
     \]
     Observe that 
    $ [\mathbf{K}_i^{I/\m I}(\textbf{x})]_j=0$ for $j>i+D$. Therefore, if 
    $z\in [\mathbf{K}_i^{P/I}(\textbf{x})]_j$
    with $j>i+D$, there is a unique homogeneous preimage $\varphi^{-1}(z)\in [\mathbf{K}_i^{P/\m I}(\textbf{x})]_j.$
    Let $\mathfrak B$ be a homogeneous basis for $ H_{\geq 1}(\mathbf{K}^{P/ I}(\textbf{x}))$ with trivial Massey operation $\mu$ on $\mathfrak B$ and let $\mathfrak B'$ be a homogeneous basis for $ H_{\geq 1}(\mathbf{K}^{P/\m I}(\textbf{x}))$. We will use $\mu$ to define a trivial Massey operation $\mu'$ on $\mathfrak B'$. First, extend $\mu$ multilinearly to $\bigsqcup_{q\geq 1}H_{\geq 1}(\mathbf{K}^{P/I}(\textbf{x}))^{q}$; note that equation (\ref{massey}) continues to hold for all $h_1, \ldots, h_q\in H_{\geq 1}(\mathbf{K}^{P/I}(\textbf{x}))$; see \cite[Lemma 11.2]{avramov2013contravariant}.
    For $h\in \mathfrak B'$, if $\varphi_*(h)\neq 0$, we can pick a representative $z_h$
    such that $\mu(\varphi_*(h))=\varphi (z_h)$, define $\mu'(h)=z_h $ and if $\varphi_*(h)=0$, we can pick a representative $z_h$ such that $z_h\in I\mathbf{K}^{P/\m I}(\textbf{x})$, define $\mu'(h)=z_h $. For $(h_1,\ldots,h_{r})\in \bigsqcup _{q\geq 2} \mathfrak B'^q$, if $\varphi_*(h_i)=0$ for some $i$, define $\mu'(h_1,\ldots,h_{r})=0$. If  $\varphi_*(h_i)\neq 0$, for all $i$, let $i_{h_t}$ and $j_{h_t}$ denote the homological and respectively the internal degree of $h_t$ for each $t$. We have  
    \[\mu(\varphi_{*} (h_1),\ldots,\varphi_{*}(h_r))\in [\mathbf K_i^{P/I}({\bf x})]_j, 
    \text{with} \]
    \[i=\sum_{t=1}^r i_{h_t}+r-1 \ \text{and}\ j=\sum_{t=1}^r j_{h_t}.\] Note that we have $j_{h_t}\ge i_{h_t}+d$
    for all $t$, and therefore $j\ge i+rd-r+1\ge i+D+1.$
     Define $ \mu'(h_1,\ldots,h_r): = \varphi^{-1} (\mu(\varphi_{*} (h_1),\ldots,\varphi_{*}(h_r))\in [\mathbf{K}_i^{P/\m I}(\textbf{x})]_j$ (as observed earlier, there is a unique such element).

    It is clear that $\mu'$ satisfies conditions $(i)$ and $(iii)$ in the definition of a trivial Massey operation. It remains to check condition $(ii)$. Let $(h_1,h_2)\in \mathfrak B'^2$. If $\varphi_*(h_1) $ or $\varphi_*(h_2)$ is zero then one of $z_{h_1}$ or $z_{h_2}$ has coefficients in $I$, and the other one has coefficients in $\m$; it follows that $\overline{z_{h_1}}\cdot z_{h_2}=0=\partial (\mu'(h_1, h_2))$. If both $\varphi_*(h_1) $ and $\varphi_*(h_2)$ are not zero, then \begin{equation*}
    \begin{aligned}
    \partial(\mu'(h_1,h_2))&=\partial(\varphi^{-1}(\mu(\varphi_*(h_1),\varphi_*(h_2))))\\
    &=\varphi^{-1}(\partial\mu(\varphi_*(h_1),\varphi_*(h_2)))\\
    &=\varphi^{-1}(\overline{\mu(\varphi_*(h_1))}\cdot \mu(\varphi_*(h_2)))\\
    &=\varphi^{-1}(\overline{\varphi(z_{h_1})}\cdot\varphi(z_{h_2}))\\
    &=\overline{z_{h_1}}\cdot z_{h_2}\\
    &=\overline{\mu'(h_1)}\mu'(h_2).  
    \end{aligned}
    \end{equation*}
    Now, let $q>2$. If $ \varphi_*(h_i)=0$ for some $i$, then  $ \sum_{j=1}^{q-1}\overline{\mu'(h_1,\cdots,h_{j})}\mu'(h_{j+1},\cdots,h_q)=0=\partial(\mu'(h_1, \cdots, h_r))$. If $ \varphi_*(h_i)$ is not zero for all $i$, then\begin{equation*}
    \begin{aligned}
    \partial(\mu'(h_1,\ldots,h_q))&=\partial(\varphi^{-1}(\mu(\varphi_*(h_1),\ldots,\varphi_*(h_q))))\\
    &=\varphi^{-1}(\partial((\mu(\varphi_*(h_1),\ldots,\varphi_*(h_q)))))\\
    &=\varphi^{-1}(  \sum_{j=1}^{q-1}\overline{\mu(\varphi_*(h_1),\cdots,\varphi_*(h_{j}))}\mu(\varphi_*(h_{j+1}),\cdots,\varphi_*(h_r)))\\
    &=  \sum_{j=1}^{q-1}\overline{\varphi^{-1}\mu(\varphi_*(h_1),\cdots,\varphi_*(h_{j}))}\varphi^{-1}\mu(\varphi_*(h_{j+1}),\cdots,\varphi_*(h_r))\\
    &=  \sum_{j=1}^{q-1}\overline{\mu'(h_1,\cdots,h_{j})}\mu'(h_{j+1},\cdots,h_r)
    \end{aligned}
    \end{equation*}
\end{proof}

\begin{corollary}
    Let $(P,\m,\kk)$ be a standard graded polynomial ring and $I\subseteq \m^2$ a proper homogeneous ideal. Let $d$ and $D$ be the lowest and respectively the highest degree of a minimal generator of $I$. If $D\leq 2(d-1)$, then $ \m^{k}I$ is Golod if and only if $k\geq\mathrm{gt}_{\m}(I)_{}$. 
\end{corollary}




% We will show that $\m^k I_{a,5}$ is Golod for all $a,k\geq 1$. Let $Z_i$ and $H_i$ be respectively the $ i$-th cycle and $i$-th homology in the Koszul complex for variables over $P/\m^k I_{a,5}$. For degree reasons we have $Z_5\cdot Z_{\geq 1}=0$, $Z_4\cdot Z_{\geq 2}=0$, and $Z_3\cdot Z_{\geq 3}=0$. Now we will show that there exists a basis for $H_1$ and $H_2$ with representative $\mathfrak R_1$ and $\mathfrak R_2$ such that $\mathfrak R_1\cdot Z_{\geq 1}=0$, and $\mathfrak R_2\cdot Z_{\geq 1}=0$. From \cite{dao2022monomial} we know that $H_1$ and $H_2$ has monomial basis. Let $[x^{\alpha}dx_1]$ be a basis element for $H_1$. If there exists $\{i,j\}\in A_5$, with $x_i^ax_j^a|x^{\alpha}$, then $x^{\alpha}dx_1\cdot Z_{i}=0 \;\forall i\geq 1$. Now, if for all $\{i,j\}\in A_5 $, $ x_i^ax_j^a$ does not divides $x^{\alpha}$. In this case we must have $\alpha_1=a-1$, $|\alpha|\geq k+2a-1\geq 2a $ and either $\alpha_3\geq a$ or $\alpha_4\geq a$. Suppose $\alpha_3\geq a$ ($\alpha_4\geq a$ is same upto change of index). For $|\alpha|\geq 2a$ to be true we must have at least one of $\alpha_2$, $\alpha_3-a$, $\alpha_4$, or $\alpha_5$ to be positive. Suppose $\alpha_2\geq 1$ (the remaining cases are similar), then \[
%         \partial(x_1^{a-1}x_2^{\alpha_2-1}x_3^{\alpha_3}x_{4}^{\alpha_4}x_5^{\alpha_5}dx_1dx_2)=-x_1^{a-1}x_2^{\alpha_2}x_3^{\alpha_3}x_{4}^{\alpha_4}x_5^{\alpha_5}dx_1+x_1^{a}x_2^{\alpha_2-1}x_3^{\alpha_3}x_{4}^{\alpha_4}x_5^{\alpha_5}
%         dx_2
%         \]
%         \[
%         \implies [(x_1^{a-1}x_{3}^a)x_2^{\alpha_2}x_3^{\alpha_3-a}x_{4}^{\alpha_4}x_5^{\alpha_5}dx_1]=[(x_1^{a}x_{3}^a)x_2^{\alpha_2-1}x_3^{\alpha_3-a}x_{4}^{\alpha_4}x_5^{\alpha_5}dx_2]\text{ in $H_1$.}
%         \]
%         Therefore, there exists a representative $h_{\alpha}$ for $[x^{\alpha}dx_1]$ such that $h_{\alpha}\cdot Z_{\geq 1}=0$. Let $x^{\alpha}dx_F$, be a monomial cycle in $Z_2$. If there exists $\{i,j\}\in A_5$, with $x_i^ax_j^a|x^{\alpha}$, then we have $|\alpha|=k+2a-1$, which implies that $x^{\alpha}dx_F\cdot Z_{\geq 1}=0$. If for all $\{i,j\}\in A_5 $, $ x_i^ax_j^a$ does not divides $x^{\alpha}$. Here up-to re indexing $F=\{1,2\}$ or $F=\{1,3\}$. Suppose $F=\{1,2\}$. In this case, we must have $\alpha_1=\alpha_2=a-1$. For $x_2x^\alpha$ to be in $\m^kI_{a,5}$, we must have $\alpha_{4}\geq a$ or $\alpha_{5}\geq a$ and for $x_1x^\alpha$ to be in $\m^kI_{a,5}$, we must have $\alpha_{3}\geq a$ or $\alpha_{4}\geq a$. But this reduces to only the following two cases \begin{itemize}
%             \item If $\alpha_5\geq a$ and $\alpha_3\geq a$, then the cycle is zero.
%             \item If $\alpha_4\geq a$, then the product of $x^{\alpha}dx_{12}$ with any monomial cycle in $Z_2$ is zero. Now, let $ z=\sum_{1\leq i_1< i_2< i_3\leq 5} f_{\{i_1,i_2,i_3\}}dx_{i_1,i_2,i_3}\in Z_3$, we have $z\cdot x^\alpha dx_{12}=\pm f_{3,4,5}x^{\alpha}dx_{12345}$. If $x_1,x_2$ divides $f_{3,4,5}$, then $z\cdot x^\alpha dx_{12}=0$. If $f_{3,4,5}=x_3^{p_3}x_4^{p_4}x_5^{p_5}$, then \[
%             z\cdot x^{\alpha}dx_{12}=\pm (x_1^{a-1}x_2^{a-1}x_4^a)x_3^{p_3+\alpha_3}x_{4}^{\alpha_4-a+p_4}x_5^{p_5+\alpha_5}dx_{12345}
%             \]\[
%             \implies \partial(z\cdot x^{\alpha}dx_{12})=\pm x_1^{a-1}x_2^{a-1}x_3^{p_3+\alpha_3+1}x_{4}^{\alpha_4+p_4}x_5^{p_5+\alpha_5}dx_{1245}\pm\]
%             \[ x_1^{a-1}x_2^{a-1}x_3^{p_3+\alpha_3}x_{4}^{\alpha_4+p_4+1}x_5^{p_5+\alpha_5}dx_{1235}\pm x_1^{a-1}x_2^{a-1}x_3^{p_3+\alpha_3}x_{4}^{\alpha_4+p_4}x_5^{p_5+\alpha_5+1}dx_{1234}
%             \]
%             Since $ z\cdot x^{\alpha}dx_{12}$ is a cycle we must have the coefficient of $dx_{1235}$ in $ \partial(z\cdot x^{\alpha}dx_{12})$, but this is only possible if $p_3+\alpha_3$ and $p_5+\alpha_5$ both are atleast $a$, but this would imply that $z\cdot x^{\alpha}dx_{12}=0$ to begin with. 
%         \end{itemize}
%         Now, Suppose $F=\{1,3\}$. In this case, we must have $\alpha_1=\alpha_3=a-1$. For $x_3x^\alpha$ to be in $\m^kI_{a,5}$, we must have $\alpha_{1}\geq a$ or $\alpha_{5}\geq a$ and for $x_1x^\alpha$ to be in $\m^kI_{a,5}$, we must have $\alpha_{3}\geq a$ or $\alpha_{4}\geq a$. Therefore we must have $\alpha_4\geq a$ and $\alpha_5\geq a$. Clearly, the product of $x^{\alpha}dx_{13}$ with any monomial cycle in $Z_2$ would be zero. Now, let $ z=\sum_{1\leq i_1< i_2< i_3\leq 5} f_{\{i_1,i_2,i_3\}}dx_{i_1,i_2,i_3}\in Z_3$, we have $z\cdot x^\alpha dx_{13}=\pm f_{2,4,5}x^{\alpha}dx_{12345}$. If $x_1,x_3$ divides $f_{2,4,5}$, then $z\cdot x^\alpha dx_{13}=0$. If $f_{2,4,5}=x_2^{p_2}x_4^{p_4}x_5^{p_5}$, then \[
%             z\cdot x^{\alpha}dx_{13}=\pm (x_1^{a-1}x_3^{a-1}x_4^ax_{5}^a)x_2^{p_2+\alpha_2}x_{4}^{\alpha_4-a+p_4}x_5^{p_5-a+\alpha_5}dx_{12345}
%             \]\[
%             \implies \partial(z\cdot x^{\alpha}dx_{13})= \pm(x_1^{a-1}x_3^{a-1})x_2^{p_2+\alpha_2+1}x_{4}^{\alpha_4+p_4}x_5^{p_5+\alpha_5}dx_{1345}\pm\]
%             \[ (x_1^{a-1}x_3^{a-1})x_2^{p_2+\alpha_2}x_{4}^{\alpha_4+p_4+1}x_5^{p_5+\alpha_5}dx_{1235}\pm(x_1^{a-1}x_3^{a-1})x_2^{p_2+\alpha_2}x_{4}^{\alpha_4+p_4}x_5^{p_5+\alpha_5+1}dx_{1235}
%             \]
%         Since $ z\cdot x^{\alpha}dx_{13}$ is a cycle we must have the coefficient of $dx_{1235}$ in $ \partial(z\cdot x^{\alpha}dx_{13})$, but this is only possible if $p_2+\alpha_2\geq a$, but this would imply that $z\cdot x^{\alpha}dx_{13}=0$ to begin with. Therefore we can choose a basis $\mathfrak B_1$ and $\mathfrak B_2$ with representatives $\mathfrak R_1$ and $\mathfrak R_2$ respectively such that, $\mathfrak R_1\cdot Z_{\geq 1}=0$ and $\mathfrak R_2\cdot Z_{\geq 1}=0$. Therefore, we can define a trivial Massey operation on the Koszul complex $\mathbf K^{P/\m^kI_{a,5}}(\textbf{x})$.
% \begin{comment}     
%     The actual Golod threshold for $I_a$ is $k_0=1$. We can check that $\m I_a$ and $\m^2 I_a$ is $d$-Golod (for definition check, \cite{herzog2018koszul}). Now, for $k\geq 3$, we have $\m^{k-1}$ is $d$-Golod, and since the product of $d$-Golod ideal is $d$-Golod, we have $\m^k I$ is $d$-Golod ideal. For $a=1$, $\m^k I$ is never $d$-Golod, so the above argument fails. For $a=1$, $\m^k I$ is also not strongly Golod for all $k\geq 1$. Using \cref{another bound} or \cref{upperbound}, we get that the Golod Threshold for $I_1$ is bounded above by 2, and computation from Macaulay2 shows that $\m I_1$ is Golod. Therefore, the Golod threshold for $\m I_1$ is also $1$. 
% \end{comment}
% \end{example}

% \begin{example}\label{cycle ideal}
% Let $S=\kk[x_1,\cdots x_{2p}]$ with $p\ge 3$, and let 
% $I=(x_i^ax_j^a\, | \, |i-j|\ne 1, 2p-1)$.
% $I$ is a Gorenstein ideal with $\mathrm{reg}(I)=?$.
% We will identify a lower bound $k_{p, a}:=(a-1)(p-2)$ for the Golod threshold of $I$.
% For $k\le k_{p, a}$, we identify two cycles $\xi_1, \xi_2$ in the Koszul complex over $R=S/\m^kI$ such that the product $\xi_1\xi_2$ is not a boundary. 
% Let
% $$
% \xi_1:=\begin{array}{c} x_1^ax_2^{a-1}x_4^{a-1}\cdots x_{2p-2}^{a-1}dx_2dx_4\cdots dx_{2p-2}+
% \\x_1^{a-1}x_2^ax_4^{a-1}\cdots x_{2p-2}^{a-1}dx_1dx_4\cdots dx_{2p-2}\\
% \end{array}
% $$
% $$
% \xi_2=\begin{array}{c}
% x_3^{a-1}x_5^{a-1}\cdots x_{2p-1}^ax_{2p}^{a-1}dx_3dx_5\cdots \hat{dx_{2p-1}} dx_{2p}+\\
% x_3^{a-1}x_5^{a-1}\cdots x_{2p-1}^{a-1}x_{2p}^a dx_3dx_5\cdots dx_{2p-1}\\
% \end{array}
% $$
% These are cycles in the Koszul complex over $R=S/\m^kI$ provided that $(a-1)(p-2)\ge k$.
% We have 
% $$
% \xi_1\xi_2=\pm x_1^a(x_2\cdots x_{2p-1})^{a-1}x_{2p}^a dx_2dx_3\cdots dx_{2p-1}
% $$
% We claim that $\xi_1\xi_2$ is not a boundary in the Koszul complex. Assume by contradiction that it is a boundary. Since the differential in the Koszul complex preserves multi-degree, the only possibility is
% $$
% xi_1\xi_2=\partial(\alpha (x_1\cdots x_{2p-1})^{a-1}x_{2p}^a dx_1\cdots dx_{2p-1} + \beta x_1^a(x_2\cdots x_{2p})^{a-1} dx_2\cdots dx_{2p})
% $$
% where $\alpha, \beta \in \kk$. However, it is easy to see that this is not the case. 

% {\em Another choice for the cycles:}
% $$
% \xi _1=(x_1\cdots x_{p-1})^{a-1} x_{p+1}^a dx_1\cdots dx_{p-1}
% $$
% $$
% \xi_2=(x_{p+2}\cdots x_{2p})^{a-1} x_p^a dx_{p+2}\cdots dx_{2p}
% $$
% These are cycles in the Koszul complex over $S/\m^kI$ for $k\le (a-1)(p-2)$, and 
% $$
% \xi_1\xi_2=(x_1\cdots x_{p-1})^{a-1}(x_px_{p+1})^a(x_{p+2}\cdots x_{2p})^{a-1}(dx_1\cdots dx_{p-1})(dx_{p+2}\cdots x_{2p})
% $$
% is not a boundary.
% Indeed, the only elements of the same multidegree that could possibly map to it are linear combinations of $$(x_1\cdots x_p)^{a-1}x_{p}^a(x_{p+2}\cdots x_{2p})^{a-1}dx_1\cdots \widehat{dx}_{p}\cdots dx_{2p}\ \mathrm{and}$$
% $$
% (x_1\cdots x_p)^{a-1}x_{p+1}^a(x_{p+2}\cdots x_{2p})^{a-1}dx_1\cdots \widehat{dx}_{p+1}\cdots dx_{2p}$$ 
% \bigskip
% \end{example}




% \begin{example}
% Let $I=(x_1^a,x_2^a,x_3^a,x_4^a)$ and $a\geq 2$. Then $\mathrm{codim}(I)=4$ and $\reg(I)=4a-3$.
% \begin{center}
% \begin{tabular}{ |c|c|c|c| }
% \hline
%  i & $t(i)-i-D+1$ & $\mathrm{reg}(I)-T(i)+i-D+1$ & $\min$ \\ 
%  \hline
%  1 & 0 & $3a-3$ & 0 \\  
% \hline
%  2 & $a-1$ & $2a-2$ & $a-1$ \\
%  \hline
%  3 & $2a-2$ & $a-1$ & $a-1$ \\  
%  \hline
%  4 & $3a-3$ & 0 & 0 \\
% \hline
% \end{tabular}
% \end{center}

% Therefore, $ \mathrm{gt}_\m(I)$ is bounded below by $a$. From $\cref{another bound}$, we have $\mathrm{gt}_\m(I)\leq a$. Therefore, $\mathrm{gt}_\m(I)=a$.   
% \end{example}
% \begin{example}
% Let $I=(x_1^2,x_2^2,x_3^2,x_4^D)$ where $D>>2$. Then $\mathrm{codim}(I)=4$, $\reg(I)=3+D$. Observe that for $i\leq 3$, $t(i)-i-D+1$ is negative and for $i=4$, $\mathrm{reg}(I)-T(i)+i-D+1$ is negative. Therefore, for this example, the lower bound $\ell$ is negative.
% \end{example}
 % \begin{obs}
 %     The lower bound in \cref{lower bound} is positive if and only if there exists $i\leq \mathrm{codim}(I)$ such that the ith shifts are squeezed between $D+i$ and $\reg(I)-D$, i.e., for all $j$ with $\beta_{ij}\neq 0$ \[
 %     i+D\leq j\leq \reg(I)-D.
 %     \]
 % \end{obs}

% In characteristic 0, for any homogeneous ideal $I$ it was shown in \cite{herzog2013ordinary} that $I^k$ is Golod for every $k\geq 2$. In the characteristic $p$ case, for any monomial ideal $I$ it was shown in \cite{herzog2018koszul} that $I^k$ is Golod for every $k\geq 2$. The main ingredient for the above results was the description of the Koszul cycles. In general, for characteristic $p$, the problem of whether $I^k$ is Golod for $k\geq 2$ is still open for general homogeneous ideals. 


% \section {Combinatorial Criterion for Monomial Koszul Basis.}
% \begin{comment}
% In this section, we will give a combinatorial criterion for a monomial ideal to have a monomial Koszul basis. To do this, we will use the Alexander dual of the Upper Koszul complex, which is also known as the Lower Koszul complex. Specifically, we will show that if the Lower Koszul complex for a monomial ideal is shellable, then the ideal has a monomial Koszul basis. Using this, we will show that ideals that look like $\sum_{i=1}^{n}x^{a_i}\m^{k_i}$ have a monomial Koszul basis. As a corollary, we will also recover the result of \cite{dao2022monomial}, which shows that monomial ideals in 3 variables have a monomial Koszul basis. \\
% \end{comment}

% For this section, let $I$ be a monomial ideal in $P=\kk[x_1,\cdots,x_n]$ and $R=P/I$. Fix $\alpha\in \mathbb{N}^n$ and let $V=\mathrm{Supp}(x^{\alpha})$. The Lower Koszul complex for $I$ in degree $\alpha$ is defined to be the simplicial complex \[\mathbf{K}_{\alpha}(I)=\{G\subseteq V:\; x^{\alpha-1_{V\setminus G}}\notin I\}.\] It is the Alexander dual of the Upper Koszul Complex, but the better way to think about $\mathbf{K}_{\alpha}(I)$ is as the complement of the upward-closed set $\mathfrak T= \{F\subseteq V: x^{\alpha-1_F}\notin I \}$. Since each element $F\in\mathfrak T$ is in one-to-one correspondence with the basis elements $x^{\alpha-1_{F}}dx_F$ for the $\kk$ vector space $[\mathbf{K}^{R}(\mathbf{x})]_{\alpha}$, the Lower Koszul complex can be used naturally to study the Koszul Homology. 

% %To study this correspondence, we will use the chain map defined in \cref{chain map} from the Koszul complex to the reduced cochain dual complex of the Lower Koszul complex.
% % \begin{setting} Let $\Delta$ be a simplicial complex. For a face $G\in \Delta$, we will use $G$ to denote the corresponding basis element in the simplicial chain complex of $\Delta$.
% % \begin{equation}\label{chain map}
% % \Phi:[\mathbf K^R(\textbf{x})]_{\alpha}\xrightarrow[]{\;\;\;\;\;\;}\widetilde{\mathbf{C}}^{\cdot}(\mathbf{K}_{\alpha}(\mathfrak a);\kk)
% % \end{equation}
% % \[
% % x^{\alpha-1_F}dx_F\xrightarrow[]{\;\;\;\;\;}(-1)^{\sigma(F)+|F|}G^*
% % \]
% % where $G=V\setminus F$, $G^*$ is the dual of $G$ and $\sigma(F)=\sum_{i\in F}\mathrm{Pos}_{V}(i)$.    
% % \end{setting}

% % \begin{proposition}
% %     $\Phi$ is a chain isomorphism.
% % \end{proposition}
% % \begin{proof}
% %     Let $F\in \mathfrak T$ and $G=V\setminus F$. The only thing we need to check is for $i\in F$, \[
% %     (-1)^{\mathrm{Pos}_{G\cup \{i\}}(i)+\sigma(F)+|F|-1}=(-1)^{\mathrm{Pos}_F(i)-1+\sigma(F\setminus \{i\})+|F\setminus \{i\}|}.
% %     \] 
% %     But this follows from the fact that $\mathrm{Pos}_{V}(i)=\mathrm{Pos}_{G\cup \{i\}}(i)+\mathrm{Pos}_{F}(i)-1$.
% %  \end{proof}   

%     % Then we have \[
%     % \Phi \circ \partial (x^{\alpha-1_F}dx_F)=\Phi\left( \sum_{\substack{i\in F\\ x^{\alpha-1_{F\setminus \{i\}}}\notin \m^kI}}(-1)^{\mathrm{Pos}_F(i)-1}x^{\alpha-1_{F\setminus \{i\}}}dx_{F\setminus \{i\}}\right)
%     % \]
%     % \[
%     % =\Phi\left( \sum_{\substack{i\in F\\ F\setminus \{i\}\in \Delta}}(-1)^{\mathrm{Pos}_F(i)-1}x^{\alpha-1_{F\setminus \{i\}}}dx_{F\setminus \{i\}}\right)
%     % \]
%     % \[
%     % = \sum_{\substack{i\in F\\ F\setminus \{i\}\in \Delta}}(-1)^{\mathrm{Pos}_F(i)-1}\Phi\left(x^{\alpha-1_{F\setminus \{i\}}}dx_{F\setminus \{i\}}\right)
%     % \]
%     % \[
%     % = \sum_{\substack{i\notin G\\ G\cup \{i\}\in \mathbf K_{\alpha}(\m^kI)}}(-1)^{\mathrm{Pos}_F(i)-1+\sigma(F\setminus \{i\})+|F\setminus \{i\}|}(G\cup \{i\})^*
%     % \]
%     % and
%     % \[
%     % \partial \circ \Phi (x^{\alpha-1_F}dx_F)=\partial \left( (-1)^{\sigma(F)+|F|}G^* \right)
%     % \]
%     % \[
%     % =(-1)^{\sigma(F)+|F|}\left(\sum_{\substack{i\notin G\\G\cup \{i\}\in \mathbf K_{\alpha}(\m^kI)}}(-1)^{\mathrm{Pos}_{G\cup \{i\}}(i)-1}(G\cup \{i\})^*\right)
%     % \]
%     % \[
%     % =\sum_{\substack{i\notin G\\G\cup \{i\}\in \mathbf K_{\alpha}(\m^kI)}}(-1)^{\mathrm{Pos}_{G\cup \{i\}}(i)+\sigma(F)+|F|-1}(G\cup \{i\})^*
%     % \]
%     % Since $\mathrm{Pos}_{\mathrm{Supp}(x^{\alpha})}(i)=\mathrm{Pos}_{G\cup \{i\}}(i)+\mathrm{Pos}_{F}(i)-1$, we have \[
%     % (-1)^{\mathrm{Pos}_{G\cup \{i\}}(i)+\sigma(F)+|F|-1}=(-1)^{\mathrm{Pos}_F(i)-1+\sigma(F\setminus \{i\})+|F\setminus \{i\}|}\]
%     % which implies $\Phi\circ \partial =\partial \circ \Phi$. Isomorphism follows as $F$ and $G$ are in one-to-one correspondence.



% % \begin{corollary}\cite{miller2005combinatorial}
% %     For $p\geq 1$\[
% %     H_{p}([\mathbf{K}^R(\textbf{x})]_{\alpha})\cong H^{|V|-p-1}(\mathbf{K}_{\alpha}(\mathfrak a);\kk).
% %     \]
% % \end{corollary}
% % The isomorphism in Proposition 3.1. gives a nice way to study the Koszul homology using the Lower Koszul Complexes. This also shows that homology for Lower Koszul complex has faces as their basis iff we have a monomial basis for the Koszul complex. In other words, if we want to show that homology for the Koszul complex has a monomial basis, we can just show that the Lower Koszul complex has homology generated by faces.

% \begin{lemma}\label{Facetcycle}
% Let $F\subseteq V$ and $G=V\setminus F$, suppose
% \(x^{\alpha-1_F}\notin I\).  Then $x^{\alpha-1_F}dx_F$ is a cycle in $[\mathbf{K}^R(\textbf{x})]_{\alpha}$ if and only if $G$ is a facet in $\mathbf K_\alpha(I)$.
% \end{lemma}

% \begin{proof}
% $\partial(x^{\alpha-1_F}dx_F)=0\iff
% \forall i\in F,\;\;x^{\alpha-1_{F\setminus\{i\}}}\in I
% \iff\forall i\notin G,\;\;
% G\cup\{i\}\notin \mathbf K_\alpha(I)\iff
% G \text{ is a facet in }\mathbf K_\alpha(I)$.
% \end{proof} 
%  In short, \cref{Facetcycle} says that all the monomial cycles in $[\mathbf{K}^{R}(\mathbf{x})]_{\alpha}$ come exactly from the facets of $\mathbf{K}_{\alpha}(I)$. Therefore, if $H_{\geq1}\left([\mathbf{K}^{R}(\mathbf{x})]_{\alpha}\right)$ has a monomial basis, then it must come from the facets of the Lower Koszul Complex. 

% %Define \[O(\mathbf{K}_{\alpha}(\mathfrak a)):=\langle [G^*]:\;G \text{ is a facet in }\mathbf{K}_{\alpha}(\mathfrak a)\rangle\]i.e., the $\kk$ vector subspace of $H^{*}(\mathbf{K}_{\alpha}(\mathfrak a);\kk)$ generated by duals of the  facets of $\mathbf{K}_{\alpha}(\mathfrak a)$ and let $\rho:O(\mathbf{K}_{\alpha}(\mathfrak a))\xrightarrow[]{}H^{*}(\mathbf{K}_{\alpha}(\mathfrak a);\kk)$ be the natural inclusion map.

% % \begin{theorem}The following are equivalent.
% % \begin{enumerate}
% % \item $H_{*}([\mathbf{K}^{R}(\textbf{x})]_{\alpha})$ has a monomial basis.
% % \item $H^*(\mathbf{K}_{\alpha}(\mathfrak a);\kk)$ has a monomial basis.
% % \item $\rho$ is surjective.
% % \end{enumerate}
% % \end{theorem}
% % \begin{proof}
% %     Equivalence of $(1)$ and $(2)$ is a direct consequence of Proposition 3.1. Equivalence of $(2)$ and $(3)$ follows from Lemma 3.3 and the discussion above.
% % \end{proof}
% % So, we can think of the cokernel of $\rho$ as being an obstruction to having a monomial basis for the Koszul homology. Define,
% % \[
% % \mathrm{Obs}([\mathbf{K}^{R}(\textbf{x})]_{\alpha}):=\mathrm{coker}(\rho)= H^{*}( \mathbf K_\alpha(\mathfrak a);\kk)/O(\mathbf K_\alpha(\mathfrak a))
% % \]
% % This object is precisely zero when the Koszul homology has a monomial basis. 



% \begin{theorem}\label{shellmonbasis}
%      If the lower Koszul complex for $I$ at all multi-degrees is shellable, then $H_{\geq 1}(\mathrm{ K}^R(\textbf{x}))$ admits a monomial basis.
% \end{theorem}
% \begin{proof}
%      For each multidegree $\alpha\in \mathbb{N}^n$, consider the chain isomorphism,
%      \begin{equation}\label{chain map}
% \Phi:[\mathbf K^R(\textbf{x})]_{\alpha}\xrightarrow[]{\;\;\;\;\;\;}\widetilde{\mathbf{C}}^{\cdot}(\mathbf{K}_{\alpha}(I);\kk)
% \end{equation}
% \[
% x^{\alpha-1_F}dx_F\xrightarrow[]{\;\;\;\;\;}(-1)^{\sigma(F)+|F|}(V\setminus F)^*
% \]
% where $(V\setminus F)^*$, is the dual of $V\setminus G$ and $\sigma(F)=\sum_{i\in F}\mathrm{Pos}_{V}(i)$. Since $\mathbf{K}_{\alpha}(I)$ is shellable, we have a basis for $H^{\geq 1}(\widetilde{\mathbf{C}}^{\cdot}(\mathbf{K}_{\alpha}(I);\kk))$, whose class representatives can be chosen to be dual to facets in $\mathbf{K}_{\alpha}(I)$ \cite[Theorem 4.3]{bjorner1996shellable}. From isomorphism $\Phi$, we get a corresponding basis for $H_{\geq 1}([\mathbf K^R(\textbf{x})]_{\alpha})$, and \cref{Facetcycle} implies that the basis is a monomial basis. Now, combining all the multi-graded bases gives a monomial basis for the Koszul Homology.
% \end{proof}
% % As a result of [\cite{gasharov1999lcm}, Lemma 2.2], to check the monomial basis for $H_{\geq 1}(\mathbf{K}^R(\textbf{x}))$, we just need to check that the lower Koszul complex for $I$ at multidegree $\alpha$ is shellable for only finitely many $\alpha$. 

% One immediate consequence of \cref{shellmonbasis} is that monomial ideal in at most 3 variables admits a monomial basis for its Koszul Homology; note that this was already known from \cite{dao2022monomial}. It turns out the converse of \cref{shellmonbasis} holds in 4 variables, but in more than 4 variables it is all wild west.

% \begin{proposition}\label{4varclassification}
%     Let $I$ be a monomial ideal in $\kk[x_1,x_2,x_3,x_4]$. The Koszul homology for $I$ admits a monomial basis if and only if all the lower Koszul complexes for I at every multidegree are shellable. 
% %     \begin{equation}\label{4shell}
% %     x^{\alpha-1_{\{1,2\}}},\;x^{\alpha-1_{\{3,4\}}}\notin I\quad \text{and}\quad
% % x^{\alpha-1_{\{1,3\}}},\;x^{\alpha-1_{\{1,4\}}},\;x^{\alpha-1_{\{2,3\}}},\;x^{\alpha-1_{\{2,4\}}}\in I        
% %     \end{equation}

% \end{proposition}
% \begin{proof}
%     Suppose there exists $\alpha=(\alpha_1,\cdots,\alpha_4)\in \mathbb{N}^4$, such that $ \mathbf{K}_{\alpha}(I)$ is not shellable. Since on 4 vertices there is a unique non-shellable simplicial complex, which is two disjoint one-dimensional facets, after relabeling we can assume $ \mathbf{K}_{\alpha}(I)=\langle \{1,2\},\{3,4\} \rangle $. Now, Consider the cycle $z=x^{\alpha-1_{\{2,3,4\}}}dx_{234}-x^{\alpha-1_{\{1,3,4\}}}dx_{134}$. The boundary of $[\mathbf K^R(\textbf{x})]_{\alpha}$ at homological degree 3 is generated by the differential of $x^{\alpha-1_{\{1,2,3,4\}}}dx_{1234}$; therefore, the class $[z]$ is non-zero, and since there are no monomial cycles in $[\mathbf K^R(\textbf{x})]_{\alpha}$ at homological degree 3, the class $[z]$ cannot be represented by a monomial basis. Therefore, the Koszul homology does not admit a monomial basis. 

%     If the Koszul homology does not have a monomial basis, then by \cref{shellmonbasis}, for some multidegree the Lower Koszul complex would not be shellable. 
% \end{proof}
% \begin{example}[Example 1.6, \cite{popescu2008monomial}]
%     Let $I=(  x_1,x_2)(x_1^p,x_2^p,x_3^p,x_4^p)$. For $\alpha=(1,1,p,p)$, the lower Koszul complex for $I$ is not shellable, so from \cref{4varclassification} we have that the Koszul homology does not have a monomial basis. 
% \end{example}

% Using \cref{4varclassification}, we can generalize Example 5.0.1 from \cite{dao2022monomial}.

% \begin{example}
%     Fix $k\geq 0$. Let $I_{k}=(x_1^{k},x_2^{k})(x_3,x_4)$. For $0\leq i\leq k$, let $ J_{i}=\m^i I_k$ then we have \[H_{3}([\mathbf{K}^R(\textbf{x})]_{(k+1,k+1,1,1)})\cong \kk[x_1^{k+1}x_2^kdx_{234}-x_1^{k}x_2^{k+1}dx_{134}].
%     \]
% \end{example}

% \begin{example}
%     For $n\geq 6$, let $I_n=\langle x_ix_j: |j-i|\notin \{1,n-1\} \rangle$. The ideal $I_n$ does not have a monomial basis for the Koszul homology at multidegree $\alpha=(1,1,0,1,1,0,\cdots,0)$. Indeed, corresponding to multidegree $\alpha$, the Lower Koszul complex is non-shellable and looks like $\langle \{1,2\},\{4,5\} \rangle$, the rest follows from arguments in Proposition 3.7.
% \end{example}

% \section{A New Class of Ideals with monomial basis.}
% In this section, we will show that an ideal of the form \[
% I=\sum_{i=1}^d x_i^{a_i}\m^{c_i},\quad\text{with}\quad 1\leq d\leq n\quad\text{and}\quad a_i+c_i\geq 1
% \]
% has a monomial basis for its Koszul Homology.
% In \cite{dao2022monomial}, H. Dao and A. D. Stefani proved that the product of any two monomial ideals in 3 variables is Golod, and the proof was based on the fact that the monomial ideals in 3 variables have a monomial Koszul basis, Lemma 2.5 and Proposition 2.6 from \cite{dao2022monomial}. They explicitly gave a proof for this. This monomial basis fact for 3 variables comes out naturally as a corollary of Theorem 3.6, as any simplicial complex on 3 vertices or less is always shellable. In $4$ variables, the only non-shellable complex on vertex set $[4]$ is $\langle\{1,2\},\{3,4\}\rangle$; therefore, the basis of the Koszul Complex of the monomial ideal in $ 4$ variables can only have monomial or binomial cycles.
% \begin{proposition}
%     Let $I$ be a monomial ideal in $\kk[x_1,x_2,x_3,x_4]$. The Koszul homology does not have a monomial basis iff, after renumbering, there exists $\alpha\in \mathbb{N}^4$ such that\[
% x^{\alpha-1_{\{1,2\}}},\;x^{\alpha-1_{\{3,4\}}}\notin \mathfrak a\quad \&\quad
% x^{\alpha-1_{\{1,3\}}},\;x^{\alpha-1_{\{1,4\}}},\;x^{\alpha-1_{\{2,3\}}},\;x^{\alpha-1_{\{2,4\}}}\in \mathfrak a
% \]
% \end{proposition} 

% In \cite{zoromski2024monomial}, it was shown that a monomial ideal with linear quotients and nice Koszul lifting property has a monomial basis for Koszul homology. Using this, they proved that some class of ideals has a monomial Koszul basis. Our approach yields a class of ideals that does not admit linear quotients. Below is one such example.
% \begin{example}
%     In 2 variable, let $I=(x_1^3,x_2^3)(x_1,x_2)$. It has a monomial Koszul basis, but it does not have a linear quotient.
% \end{example}
% Let $1\leq d\leq n$. For all $i\in [d]$, let $a_i,c_i\in \mathbb{N}$ such that $a_i+c_i\geq 1$. Define \[
% I=\sum_{i=1}^dx_i^{a_1}\m^{c_i}
% \]
% Fix $\alpha\in \mathbb{N}^n$, and let $V=\mathrm{Supp}(x^{\alpha})$. For all $i\in [d]$, let $\ell_i:=a_i+c_i+|V|-|\alpha|$ and define \[
% \lambda_0=\min\{\ell_i:\alpha_i>a_i\}
% \]
% and, for $i\in V$, 
% \[
% \lambda_i=\min\{\ell_i:\alpha_i=a_i\}
% \]
% with every empty minimum interpreted as $\infty$.

% \begin{proposition}The following is the combinatorial description of the lower Koszul complex of $I$ at degree $\alpha$:\[ 
%      K_{\alpha}(I)=\{ G\subseteq V:\;|G|< \lambda_0\text{ and }|G|< \lambda_i\text{ for every }i\in G\}.
%     \]
% \end{proposition}
% \begin{proof}
%     ($\subseteq$) Let $G\in \mathbf{K}_{\alpha}(I)$, then we have $x^{\alpha-1_{V\setminus G}}\notin I$. If $|G|\geq\lambda_0$, then there exists $i\in [N]$ such that $|G|\geq\ell_i$ and $\alpha_i>a_i$ which implies that $x^{\alpha-1_{V\setminus G}}\in x^{a_i}\m^{c_i}\subseteq I$, a contradiction. If for some $i\in G$, $|G|\geq \lambda_i$ then we have $i\in G$ such that $|G|\geq \ell_i$ and $\alpha_i=a_i$ which implies that $x^{\alpha-1_{V\setminus G}}\in x^{a_i}\m^{c_i}\subseteq I$, a contradiction.

%     ($\supseteq$) Let $G\subseteq V$ with $|G|<\lambda_0$ and $|G|<\lambda_i$ for every $i\in G$. If $x^{\alpha-1_{V\setminus G}}\in I$ then $x^{\alpha-1_{V\setminus G}}\in x^{a_i}\m^{c_i}$ for some $i\in [N]$. If $\alpha_i>a_i$, then $|\alpha|-|V|+|G|-a_i\geq c_i$, which implies that $|G|\geq \ell_i$ a contradiction. If $\alpha_i=a_i$, then we have $i\in G$ and $|\alpha|-|V|+|G|-a_i\geq c_i$, which implies that $|G|\geq \lambda_i$, a contradiction.
% \end{proof}

% \begin{proposition}
%     $\mathbf K_{\alpha}(I)$ is a shifted simplicial complex.
% \end{proposition}
% \begin{proof}
%     Let $N=|V|$. We label the nodes of $V$ as \[
%     V=\{v_1<\cdots <v_N\}
%     \]
%     where we have \[
%     \lambda_{v_1}\geq \cdots \geq \lambda_{v_N}
%     \]
%     Now let $G\in \mathbf K_{\alpha}(I)$. Suppose $v_i\in G$ and $j< i$. We have \[
%     |(G\setminus \{v_i\})\cup\{v_j\}|=|G|<\lambda_0
%     \]
%     and since $\lambda_{v_j}\geq \lambda_{v_i}$ we have\[
%     |(G\setminus \{v_i\})\cup\{v_j\}|=|G|<\lambda_{v_i}\leq \lambda_{v_j}
%     \]
%     which implies that $(G\setminus {v_i})\cup\{v_j\}\in \mathbf K_{\alpha}(I)$. Therefore $\mathbf K_{\alpha}(I)$ is a shifted simplicial complex.
% \end{proof}

% \begin{corollary}\label{maincorr}
%     Monomial ideals of the form $\sum_{i=1}^d x_i^{a_i}\m^{c_i}$ have monomial Koszul basis. 
% \end{corollary}
% \begin{proof}
%     From \cite{bjorner1997shellable} we know that all shifted complexes are shellable. The rest follows from \cref{shellmonbasis} we have the result.
% \end{proof}

\section{Combinatorial Criterion for Monomial Koszul Basis}
% In this section we give a combinatorial criterion for the existence of a monomial basis of Koszul homology of a monomial quotient ring. Using this criterion, we will give a sufficient and necessary condition for monomial ideals in 4 variables to have a monomial basis for their Koszul Homology. At the end end we will show that the ideal of the form $ \m^k (x_1^{a_1},\cdots,x_n^{a_n})$ admits a monomial basis for its Koszul homology, and we will use this result in next section to establish the exact formula for the Golod threshold for a monomial complete intersection.

Let $I$ be a monomial ideal in $P=\kk[x_1,\cdots,x_n]$ and $R=P/I$. Fix $\alpha\in \mathbb{N}^n$ and let $V=\mathrm{Supp}(x^{\alpha})$. The Lower Koszul complex for $I$ in degree $\alpha$ is  the simplicial complex \[\mathbf{K}_{\alpha}(I)=\{G\subseteq V:\; x^{\alpha-1_{V\setminus G}}\notin I\}.\]  $\mathbf{K}_{\alpha}(I)$ consists of the complements of the upward-closed set $\mathfrak T_{\alpha}= \{F\subseteq V: x^{\alpha-1_F}\notin I \}$, and therefore it is a simplicial complex. Note that the  elements $F\in\mathfrak T_{\alpha}$ are in one-to-one correspondence with the basis elements $x^{\alpha-1_{F}}dx_F$ of the $\kk$ vector space $[\mathbf{K}^{R}(\mathbf{x})]_{\alpha}$.
\begin{definition}
A {\em monomial cycle} in $\mathbf{K}^{R}(\mathbf{x})$ is a cycle of the form $[x^{\alpha-1_F}dx_F]$
for some $F\subseteq [n]$. 
We say that the Koszul homology of $R$ admits a {\em monomial Koszul basis} if there exists a basis of $H_{\ge 1}(\mathbf{K}^{R}(\mathbf{x}))$ consisting of monomial cycles. 
    
\end{definition}

We will now give a sufficient condition for the Koszul homology of $R$ to admit a monomial Koszul basis, which will be in terms of shellability of the lower Koszul complexes. For the definition of shellability, we refer the reader to \cite{bjorner1996shellable}.

% The following Lemma shows that the monomial cycles in $[\mathbf{K}^{R}(\mathbf{x})]_{\alpha}$ come  from the facets of $\mathbf{K}_{\alpha}(I)$.  

% \begin{lemma}\label{Facetcycle}
% Let $F\subseteq V$ and $G=V\setminus F$, suppose
% \(x^{\alpha-1_F}\notin I\).  Then $x^{\alpha-1_F}dx_F$ is a cycle in $[\mathbf{K}^R(\textbf{x})]_{\alpha}$ if and only if $G$ is a facet in $\mathbf K_\alpha(I)$.
% \end{lemma}

% \begin{proof}
% $\partial(x^{\alpha-1_F}dx_F)=0\iff
% \forall i\in F,\;\;x^{\alpha-1_{F\setminus\{i\}}}\in I
% \iff\forall i\notin G,\;\;
% G\cup\{i\}\notin \mathbf K_\alpha(I)\iff
% G \text{ is a facet in }\mathbf K_\alpha(I)$.
% \end{proof} 

 

\begin{theorem}
    \label{shellmonbasis}
     If $K_\alpha(I)$ is shellable for all multidegrees $\alpha$,  then $H_{\geq 1}(\mathrm{ K}^R(\textbf{x}))$ admits a monomial basis.
\end{theorem}
\begin{proof}
     For each multidegree $\alpha\in \mathbb{N}^n$, consider the chain isomorphism,
     \begin{equation}\label{chain map}
\Psi:[\mathbf K^R(\textbf{x})]_{\alpha}\xrightarrow[]{\;\;\;\;\;\;}\widetilde{\mathbf{C}}^{\cdot}(\mathbf{K}_{\alpha}(I);\kk)
\end{equation}
\[
x^{\alpha-1_F}dx_F\xrightarrow[]{\;\;\;\;\;}(-1)^{\sigma(F)+|F|}(V\setminus F)^*
\]
where $\sigma(F)=\sum_{i\in F}\mathrm{pos}_{V}(i)$. Since $\mathbf{K}_{\alpha}(I)$ is shellable, a basis for $H^{\geq 1}(\widetilde{\mathbf{C}}^{\cdot}(\mathbf{K}_{\alpha}(I);\kk))$ exists, whose class representatives can be chosen to be dual to facets in $\mathbf{K}_{\alpha}(I)$ by \cite[Theorem 4.3]{bjorner1996shellable}. From the isomorphism $\Psi$, we get a corresponding basis for $H_{\geq 1}([\mathbf K^R(\textbf{x})]_{\alpha})$, which is monomial by construction. The union of all the multigraded bases gives a monomial basis for $H_{\ge 1}(\mathbf K^R(\textbf{x})).$
\end{proof}
% As a result of [\cite{gasharov1999lcm}, Lemma 2.2], to check the monomial basis for $H_{\geq 1}(\mathbf{K}^R(\textbf{x}))$, we just need to check that the lower Koszul complex for $I$ at multidegree $\alpha$ is shellable for only finitely many $\alpha$. 

One immediate consequence of \cref{shellmonbasis} is that a monomial ideal in at most 3 variables admits a monomial basis for its Koszul Homology; note that this was already known from \cite{dao2022monomial}. We note that the converse of \cref{shellmonbasis} holds in 4 variables.

\begin{proposition}\label{4varclassification}
    Let $I$ be a monomial ideal in $\kk[x_1,x_2,x_3,x_4]$. The Koszul homology for $I$ admits a monomial basis if and only if all the lower Koszul complexes for I at every multidegree are shellable. 
 %    \begin{equation}\label{4shell}
 %    x^{\alpha-1_{\{1,2\}}},\;x^{\alpha-1_{\{3,4\}}}\notin I\quad \text{and}\quad
 % x^{\alpha-1_{\{1,3\}}},\;x^{\alpha-1_{\{1,4\}}},\;x^{\alpha-1_{\{2,3\}}},\;x^{\alpha-1_{\{2,4\}}}\in I        
 %     \end{equation}
 \end{proposition}
\begin{proof}
     Suppose there exists $\alpha=(\alpha_1,\cdots,\alpha_4)\in \mathbb{N}^4$, such that $ \mathbf{K}_{\alpha}(I)$ is not shellable. Up to labeling, there is a unique non-shellable simplicial complex on 4 vertices, which is two disjoint one-dimensional facets. Therefore we can assume $ \mathbf{K}_{\alpha}(I)=\{ \{1,2\},\{3,4\} \} $.  Consider the cycle $z=x^{\alpha-1_{\{2,3,4\}}}dx_{234}-x^{\alpha-1_{\{1,3,4\}}}dx_{134}$. The boundary of $[\mathbf K^R(\textbf{x})]_{\alpha}$ at homological degree 3 is generated by the differential of $x^{\alpha-1_{\{1,2,3,4\}}}dx_{1234}$; therefore, the class $[z]$ is non-zero, and since there are no monomial cycles in $[\mathbf K^R(\textbf{x})]_{\alpha}$ at homological degree 3, the class $[z]$ cannot be represented by a monomial basis. Therefore, the Koszul homology does not admit a monomial basis. 
The other implication is known from \cref{shellmonbasis}. 
 \end{proof}

% \begin{example}
%     Fix an integer $k\geq 1$. Let $I_{k}=(x_1^{k},x_2^{k})(x_3,x_4)$. For $\alpha=(k,k,1,1)$, the lower Koszul complex for $I_{k}$ is not shellable; therefore, the Koszul homology for $I_k$ does not have a monomial basis. As one can check, \[H_{3}([\mathbf{K}^R(\textbf{x})]_{(k,k,1,1)})\cong \kk[x_1^kx_2^{k-1}dx_{234}-x_1^{k-1}x_2^{k}dx_{134}].
%     \]
% \end{example}

Throughout the rest of this section, let $P=k[x_1,\ldots,x_n]$ be a standard multigraded polynomial ring with homogeneous maximal ideal $\mathfrak m=(x_1,\ldots,x_n)$, and let
\[
I=(x_1^{a_1},\ldots,x_n^{a_n}), \quad a_i\geq 1 \text{ for all } 1\leq i\leq n,
\]
be a monomial complete intersection. Set $R=P/I$. We will use \cref{shellmonbasis} to show that the Koszul homology of $\m^k I$ admits a monomial basis for all $k\geq 0$.

Fix a multidegree $\alpha\in \mathbb{N}^n$. For each $i\in [n]$, let $q_i:=|\alpha| -a_i-k$ and $q'_i:=|\mathrm{Supp}(x^{\alpha}) | - q_i$. Recall that the lower Koszul complex $\mathbf{K}_{\alpha}(\m^kI)$ is the complement of the upward-closed set \[
\mathfrak T_{\alpha}=\{F\subseteq \mathrm{Supp}(x^{\alpha}): x^{\alpha-1_F}\notin \m^kI \}
\]

\begin{obs}\label{obs1} Let $F\subseteq \mathrm{Supp}(x^{\alpha})$. We have $F\in \mathfrak T_{\alpha}$ if and only if for all $i\in [n]$ with $\alpha_i\ge a_i$, we have ($\alpha_i=a_i$ and $i\in F$) or $|F|>q_i$.

\end{obs} 

\begin{obs}\label{obs2}
For $G\subseteq \mathrm{Supp}(x^{\alpha})$, we have $G\in \mathbf K_{\alpha}(I)$ if and only if for all $i$ with $\alpha_i\ge a_i$, we have ($\alpha_i=a_i$ and $i\notin G$) or $|G|<q_i'$. Equivalently, $G\in \mathbf K_{\alpha}(I)$ if and only if the following hold:

\begin{enumerate}
\item $|G|<\mathrm{min}\{q_i': i\in [n]\ \mathrm{with}\ \alpha_i>a_i\}$, and 

\item $|G|<\mathrm{min}\{q_i': i\in G \ \mathrm{with}\ \alpha_i=a_i\}$

\end{enumerate}

\end{obs}
\begin{proof}
Apply Observation \ref{obs1} to $F:=\mathrm{Supp}(x^{\alpha})\, \backslash\, G$.
\end{proof}

\begin{proposition}\label{shiftedsimplex}
There is a suitable ordering of the variables $x_1, \ldots, x_n$ such that the simplicial complex $\mathbf K_{\alpha}(I)$ is shifted, i.e. if $G\in \mathbf K_{\alpha}(I)$, $j\in G$ and $j'\notin G$ such that $j'<j$, then $G':=(G\, \backslash \{j\}) \cup \{j'\} \in \mathbf K_{\alpha}$.
\end{proposition}
\begin{proof}
Reorder the variables $x_1, \ldots, x_n$ in such a way that the indices $i$ with $\alpha_i\ne a_i$ (if any) are in the beginning of the list, i.e. 
there exists $s\in [n]$ such that $\alpha_i=a_i \Leftrightarrow i\ge s+1$, and $q'_{s+1}\ge q'_{s+2}\ge \ldots \ge q'_n$.

Let $G\in \mathbf K_{\alpha}(I)$, $j\in G$ and$j'\notin G$ with $j'<j$, and let $G':=(G\, \backslash \{j\}) \cup \{j'\}$. We need to prove that $G'$ satisfies conditions (1) and (2) from Observation (\ref{obs2}). Note that $|G'|=|G|$, so condition (1) for $G'$ follows immediately from condition (1) for $G$. The ordering of the variables together with the assumption $j'<j$ implies that 
$$
\mathrm{min}\{q_i': i \in G \ \mathrm{with} \ \alpha_i=a_i\} \le \mathrm{min}\{q_i': i \in G'\  \mathrm{with} \ \alpha_i=a_i\},
$$
and condition (2) for $G'$ also follows from the corresponding condition for $G$.
\end{proof}
\begin{corollary}\label{monomialbasisforCI}
    The Koszul homology of  $\m^k I$ admits a monomial basis for all $k\geq 0$.
\end{corollary}
\begin{proof}
    From \cref{shiftedsimplex} we know that the lower Koszul complex for the ideal $\m^k I$ is shifted at all multidegrees. Since all  shifted complexes are shellable \cite[Theorem 11.3]{bjorner1997shellable}, \cref{shellmonbasis} implies that the Koszul homology of $\m^k I$ admits a monomial basis.
\end{proof}

\begin{remark}
    Using a similar argument, one can also show that the conclusion of \cref{monomialbasisforCI} also holds for ideals of the form $\sum_{i=1}^nx_i^{a_i}\m^{c_i}$. 
\end{remark}


\section{The Golod Threshold for monomial complete intersection ideals}
% In this section, we will compute the Golod threshold for a monomial complete intersection. 
Throughout this section, let $P=k[x_1,\ldots,x_n]$ be a standard graded polynomial ring with homogeneous maximal ideal $\mathfrak m=(x_1,\ldots,x_n)$, and let
\[
I=(x_1^{a_1},\ldots,x_n^{a_n}), \quad a_i\geq 1 \text{ for all } 1\leq i\leq n,
\]
be a monomial complete intersection. Set $R=P/I$. 

 For $F\subseteq [n]$ with $|F|\geq 2$, let \[
a_F:=\sum_{i\in K} a_i-|F|=\sum_{i\in F}(a_i-1)
\]
and \[
m_F:=\mathrm{min}\{a_{F'}\, | \, F'\subseteq F, |F'|=|F|-1\}
\]
For a partition $K\sqcup L= [n]$ with both $K$ and $L$ consisting of at least two elements, let $m_{K, L}:=\mathrm{min}\{m_K, m_L\}$.
(The notation $K\sqcup L=[n]$ means that $K\cup L=[n]$ and $K\cap L=\emptyset$.)
 \begin{theorem}\label{preciseCIbound}
     We have 
    $$
    \mathrm{gt}_{\m}(I)=1+\mathrm{max}\{m_{K, L}: K\sqcup L= [n]\;\text{and}\;|K|,|L|\geq 2 \}.
     $$
 \end{theorem}

% \begin{lemma}\label{notgolod}
%  If $K\sqcup L$ is a partition of $\{1\, \ldots, n\}$ with $|K|, |L|\ge 2$, then $\m^{m_{K, L}}I$ is not Golod.
% \end{lemma}
% \begin{proof}
% We have $$\prod_{i\in K} x_i^{a_i-1} \in (\m^{m_K}I):(x_i \, | \, i\in K)\subseteq (\m^{m_{K,L}}I):(x_i \, | \, i\in K)$$ and 
% $$\prod_{i\in L} x_i^{a_i-1}\in (\m^{m_L}I):(x_i\, |\, i \in L)\subseteq (\m^{m_{K,L}}I):(x_i \, | \, i\in L),$$  but $\prod_{i=1}^n x_i^{a_i-1}\notin \m^{m_{K,L}}I$. 
% \end{proof}
% This is the same argument used by Stefani in \cite{de2016products}, to show that the product of monomial ideals in 4 variables may not be Golod.

 We will need the following description of a basis of $H_{\geq 1}(\mathbf{K}^R(\mathbf{x}))$.  

\begin{proposition}\label{basischoice}
     There is a basis of $H_{\geq 1}(\mathbf{K}^R(\mathbf{x}))$ consisting elements of the form  $[udx_F]$ where $F\subseteq [n]$ is nonempty and $u\in m^{k-1}I$
     or $u=\prod_{i\in F} x_i^{a_i-1}$.        
\end{proposition}
\begin{proof}
    % Let $\Theta$ be the $k$-linear subspace of $H_{\geq 1}(\mathbf{K}^R(\mathbf{x}))$ generated by the elements whose representatives are of the above form. 
    
    We know from \cref{monomialbasisforCI}  that $H_{\geq 1}(\mathbf{K}^R(\mathbf{x}))$ has a basis consisting of elements of the form $[udx_{F}]$ where $u$ is a monomial and $F\subseteq [n]$ is nonempty. It remains to show that we can choose the monomials $u$ as in the statement. Since $udx_F$ is a cycle in the Koszul complex, we have $ux_i\in \m^kI$ for all $i\in F$. 
    This implies that $u\in \m^{k-1}I$ or $u$ is a multiple of $\prod_{i\in F}x_i^{a_i-1}$.
    Indeed, for each $i$ there exists $q_i$ such that $ux_i=hx_{q_i}^{a_i}$ with $h\in \m^k$. If $q_i\ne i$ then $x_i$ divides $h$ and it follows that $u\in \m^{k-1}I$. If $q_i=i$ for all $i\in F$, it follows that $u$ is a multiple of $\prod_{i\in F}x_i^{a_i-1}.$
    Suppose $u=M\prod_{i\in F} x_i^{a_i-1}$ with $M$ a monomial. If $M$ is a unit, our element has the required form.  Otherwise, pick $\ell \in [n]$ such that $x_{\ell}$ divides $M$. If $\ell \in F$, then $u=x_{\ell}^{a_\ell} M'\prod_{i\ne \ell, i\in F} x_i^{a_i-1}$ where $M'=M/x_{\ell}$, and it follows that $u\in \m^{k-1}I$. If $\ell \notin F$, 
    consider $y:=(u/x_{\ell})dx_{F\cup \{\ell\}}$. We have 
    \[
    \partial(y)=\pm udx_F+\sum_{i\in F}\pm (x_iu/x_{\ell})dx_{F\cup \{\ell\}\setminus \{i\}} 
    \]
    Note that $x_iu/x_{\ell} \in \m^{k-1}I$ for all $i\in F$, and therefore $[udx_F]$ is a linear combination of elements of the form $[u'dx_{F'}]$ with $u'\in \m^{k-1}I$ and $F'\subseteq [n]$.
    
    % Here, $ x_iu/x_l \in \mathfrak m^{k-1}I$. Therefore, we have \[
    % [udx_F]=\sum_{i\in F}\pm [(x_iu/x_l)dx_{F\cup \{l\}\setminus \{i\}}]\in \Theta
    % \]

    
    
    
    % to prove the above claim it is enough to show that an element in $H_{\geq 1}(\mathbf{K}^R(\mathbf{x}))$ represented by a monomial cycle is in the subspace $\Theta$. Let $udx_F$ be a monomial cycle in the Koszul complex and assume that $u\notin \mathfrak m^{k-1}I$. Now, as $udx_F$ is a cycle we have $u\in (\mathfrak m^k I:(x_i|i\in F))$. So, for each $i\in F$, there exists $q_i\in [n]$ such that \[
    % x_{q_i}^{a_{q_i}}|x_iu\quad \text{and}\quad \deg(x_iu/x_{q_i}^{a_{q_i}})\geq k 
    % \]
    % Now, if $q_i\neq i$, then \[\deg(x_iu/x_{q_i}^{a_{q_i}})\geq k\implies \deg(u/x_{q_i}^{a_{q_i}})\geq k-1\implies u\in \mathfrak m^{k-1}I\]
    % which is a contradiction. Therefore, $x_i^{a_i-1}|u$ for all $i\in F$. Now, if $x_i^{a_i}|u$ for some choice of $i\in F$, then $u\in \mathfrak m^{k-1}I$, which is again a contradiction. 
    % Therefore \[
    % u=M\prod_{i\in F}x_i^{a_i-1}\quad \text{where}\quad x_i\not{|}M\;\;\forall i\in F.
    % \]
    % If $M$ is a unit, then we are done, 
     
\end{proof}

\begin{lemma}\label{deg cond}
  Let $F\subseteq [n]$. The element $\prod_{i\in F} x_i^{a_i-1}dx_F$ is a cycle in $\mathbb K^R({\bf x})$ if and only if $m_F\ge k$.  
\end{lemma}
\begin{proof}
$\prod_{i\in F} x_i^{a_i-1}dx_F$ is a cycle in $\mathbb K^R({\bf x})$ if and only if for all $i\in F$ we have  \[\displaystyle x_i^a\prod_{j\in F\, \backslash \{i\}} x_j^{a_j-1} \in \m^kI,  \ \text{which is equivalent to} \prod_{ j\in F\, \backslash \, \{i\}} x_j^{a_j-1} \in \m^k\] 
This is clearly equivalent to $m_F\ge k$.
\end{proof}



We are now ready to prove \cref{preciseCIbound}
\begin{proof} 
Let $k>\mathrm{max}\{m_{K, L}: K\sqcup L= [n]\;\text{and}\;|K|,|L|\geq 2 \}$. We will use \cref{zeroproduct} to prove that $\m^kI$ is Golod.
More precisely, let $\mathbf{B}$ be a basis of $H_{\geq 1}(\mathbf{K}^R(\mathbf{x}))$ as in \cref{basischoice}. 
We claim that for every two $[z]=[udx_F]$ and $[z']=[u'dx_K]$ in $\mathbf B$, we have 
$z\cdot z'=0$

From \cref{basischoice}, we have $u\in \m^{k-1}I$ or $u=\prod_{i\in F} x_i^{a_i-1}$, and   $u'\in \m^{k-1}I$ or $u'=\prod_{i\in K} x_i^{a_i-1}$.
If $u\in \m^{k-1}I$ or $u'\in \m^{k-1}I$, then $uu'\in \m^kI$ and therefore $z\cdot z'=0$. 

Assume $u=\prod_{i\in F} x_i^{a_i-1}$ and $u'=\prod_{i\in K} x_i^{a_i-1}$. Note that we must have $|F|, |K|\ge 2$. 
Assume $F\cap K=\emptyset$.
From \cref{deg cond} we have $m_F\ge k$ and $m_K\ge k$.
Let $L:=[n]\, \backslash\,  F$. Since $F\subseteq L$, we have $m_L\ge m_F$, and therefore $m_L\ge k$.
It follows that $m_{K, L} \ge k$. This contradicts the choice of $k$, therefore $F\cap K$ must be nonempty. This implies $z\cdot z'=0$.



Now let $k\le \mathrm{max}\{m_{K, L}: K\sqcup L= [n]\;\text{and}\;|K|,|L|\geq 2 \}$. 
Let $K\sqcup L =[n]$ be a partition with $|K|,\;|L|\geq 2$ such that $k\leq m_{K,L}$. \cref{deg cond} shows that $ \left[\prod_{i\in K} x_i^{a_i-1}dx_{K}\right]$ and $\left[\prod_{i\in L} x_i^{a_i-1}dx_{L}\right]$ are elements of $H_{\geq 1}(\mathbf{K}^R(\mathbf{x})) $. Note that their product $\left[\prod_{i=1}^n x_i^{a_i-1}dx_{[n]}\right]$ is not zero in $H_{\geq 1}(\mathbf{K}^R(\mathbf{x}))$. Therefore, $ \m^kI$ is not Golod. 
\end{proof}  


\section{Products of Ideals are Golod in dimension 3.}
% \begin{proposition}
%     In $P=\kk[x_1,\cdots,x_4]$, the product of any equigenerated monomial Golod ideal and the irrelevant maximal ideal is Golod. 
% \end{proposition}
% In this section, we will show that in a regular local ring of embedding dimension 3, the product of any two proper ideals is Golod.

Let $(P,\m,\kk)$ be a regular local ring of dimension 3 and let $I\subseteq \m^2$ be a proper ideal. Suppose the product structure on the positive Koszul homology over $P/I$ is trivial. Let $\mathfrak B$ be a basis for $H_{\geq 1}(\mathbf K^{P/I}(\textbf{x}))$. For $h\in \mathfrak B$, let $z_h$ be a representative for $h$. Let $\mu:\sqcup_{q\geq 1} \mathfrak B^q\xrightarrow[]{} \mathbf{K}^{P/I}(\textbf{x})$ be a map which is defined as follows: for $(h_1,\cdots,h_q)\in \mathfrak B^q$, if $q=1$, then $\mu(h_1)=z_{h_1}$, if $q=2$, then $\mu(h_1,h_2)= z$, where $\partial(z)=\overline{z_{h_1}}\cdot z_{h_2}$, if $q\geq 3$, then $\mu(h_1\ldots,h_q)=0$. One can check that $\mu$ is a trivial Massey operation. 
Therefore, a proper ideal $I\subseteq \m^2$ is Golod if and only if the product structure on the positive Koszul homology over $P/I$ is trivial, i.e., $H_{\geq 1}(\mathbf K^{P/I}(\textbf{x}))\cdot H_{\geq 1}(\mathbf K^{P/I}(\textbf{x}))=0$; see \cite{golod1962homologies}. This fact is not true in general in higher dimensions; see \cite[Theorem 3.1]{katthan2017non}.

% Let $I$ and $J$ be a proper ideals in $P$. Observe that a sufficient condition for $IJ$ to be Golod is the vanishing of one of the following maps:
%     \begin{equation}\label{golodCriterion}
%     H_{2}(\mathbf{K}^{P/IJ}(\mathbf{x}))\xrightarrow[]{}H_{2}(\mathbf{K}^{P/I}(\mathbf{x}))\quad\text{and}\quad H_{2}(\mathbf{K}^{P/IJ}(\mathbf{x}))\xrightarrow[]{}H_{2}(\mathbf{K}^{P/J}(\mathbf{x}))  
%     \end{equation}
% which are induced from the canonical projections $P/IJ\xrightarrow[]{}P/I$ and $P/IJ\xrightarrow[]{}P/J$.
%  Indeed, the vanishing of either of the above two maps allows us to choose a basis for $H_{2}(\mathbf{K}^{P/IJ}(\mathbf{x}))$ whose coefficients can be chosen to be elements of $I/IJ$ or $J/IJ$. Since a basis for $ H_1(K^{P/IJ}(\textbf{x}))$ can always be chosen to be represented by cycles with coefficients from $I$ or $J$, we have $H_{2}(\mathbf{K}^{P/IJ}(\mathbf{x}))\cdot H_{1}(\mathbf{K}^{P/IJ}(\mathbf{x}))=0$. Now, this with \cite[3.1]{vandebogert2022products} implies that $IJ$ is Golod. 

% If the ideal $I$ and $J$ are both grade 3 Gorenstein ideals in $P$, using the data of their minimal free resolution\cite{buchsbaum1977algebra} we can show that one of the maps in \cref{golodCriterion} vanishes. In general, we do not have the data of the minimal free resolution for an arbitrary ideal, but one can bypass this by embedding the ideals $I$ and $J$ into grade 3 Gorenstein ideals.


\begin{lemma}\label{mainlemma1}
    Let $(P,\m,k)$ be a regular local ring, and let $I$ and $J$ be two proper ideals in $P$. Let $\mathfrak a$ be a grade 3 Gorenstein ideal with $I\subseteq \mathfrak a$. If \begin{equation}\label{torGor}
     \Tor_{2}^P(P/IJ,\kk)\xrightarrow[]{} \Tor_{2}^P(P/\mathfrak a,\kk)   
    \end{equation}
     
    is a non-zero map, then there exists $f\in \mathfrak a$ such that $\mathfrak a\subseteq (f)+\m J$.
\end{lemma}
\begin{proof}
    Let $(G_{\cdot},\partial^{IJ}_{\cdot})$ be a minimal free resolution for $P/IJ$ over $P$, and $(F_{\cdot},\partial^{\mathfrak a}_{\cdot})$ be a minimal free resolution for $P/\mathfrak a$ over $P$. Due to the structure theorem of Buchsbaum-Eisenbud\cite{buchsbaum1977algebra} for grade 3 Gorenstein ideals, we can choose $\partial^{\mathfrak a}_2$ to be multiplication by an odd-sized alternating matrix $\Phi$, and the submaximal Pfaffians of $\Phi$ generate $\mathfrak a$. Lift the canonical projection maps $P/IJ\xrightarrow[]{}P/ \mathfrak a$ to the chain map $\alpha_{\cdot}:G_{\cdot}\xrightarrow[]{}F_{\cdot}$. Since $I\subseteq \mathfrak a$, the generators of $IJ$ are linear combinations of the generators of $\mathfrak a$ with coefficients in $J$, and therefore $\mathrm{im}(\alpha_1)\subseteq JF_1$.

Suppose \cref{torGor} is a non-zero map. This implies that $\alpha_2$ has at least one unit entry; suppose this entry is in the $ r$-th column $C_r$. Choose an $n\times n$ invertible matrix $M$ with entries in $P$ and the first column equal to $C_r$. Set $\Phi_0:=M^{T}\Phi M$, where $M^T$ is transpose of $M$. Since $M$ is invertible, the submaximal Pfaffians of both $\Phi_0$ and $\Phi$ generate the same ideal, which is $\mathfrak a$. Let $\mathbf{e}_1$ be the first canonical generator for $F_1$ and $\mathbf{b}_r$ be the $ r$-th canonical generator of $G_2$. Applying $\Phi_0$ to $\mathbf{e}_1$, we get $
\Phi_0(\mathbf e_1)= M^T \Phi M (\mathbf{e}_1)=M^T \Phi (C_r)$, and since $ \alpha_2(\textbf{b}_r)=C_r$ and $\Phi\alpha_2=\alpha_1 \partial_2$, we have 
\[
\Phi_0(\textbf{e}_1)=M^T\Phi (C_r)= M^T\Phi \alpha_2(\mathbf{b}_r)=M^T\alpha_1\partial_2(\mathbf b_r).
\] 
Since we have $\mathrm{im} (\partial_2)\subseteq \m G_1$ and $\mathrm{im}(\alpha_1)\subseteq JF_1$, we obtain \[\Phi_0(\mathbf e_1)=M^T\alpha_1\partial_2(\mathbf b_r)\subseteq M^T\alpha_1(\m G_1)\subseteq \m JF_1\]
Since $\Phi_0$ is an alternating matrix, it follows that the first row and the first column of $\Phi_0$ have entries in $\m J$. Let $f$ be the Pfaffian of $\Phi_0$ obtained by deleting the first row and column, which is one of the minimal generators of $\mathfrak a$. All the other submaximal Pfaffians are in $\m J$, and therefore $\mathfrak a\subseteq (f)+\m J$.
\end{proof}

\begin{lemma}\label{GorensteinEmbed}
    Let $I$ be a proper ideal in a Noetherian local ring $(P,\m,\kk)$. Then for any $x\notin I$, there exists a $\m$-primary Gorenstein ideal $\mathfrak a$ such that $I\subseteq \mathfrak a$ and $x\notin \mathfrak a$. 
\end{lemma}
\begin{proof}
    By the Krull intersection theorem, we have $\bigcap_{n\geq 1} (I+\m^n)=I$. Since $x\notin I$, there exists $n\geq 1$ such that $ x\notin I+\m^n$. Let $\mathfrak a$ be an irreducible component of $I+\m^n$ with $x\notin \mathfrak a$. The ideal $\mathfrak{a}$ is Gorenstein by \cite[Theorem 3.2.10]{bruns1998cohen}.
     % To show $\mathfrak a$ is Gorenstein it is enough to show that $\dim_{\kk} \mathrm{Soc}(P/\mathfrak a)=1$. If $ \dim_{\kk} \mathrm{Soc}(P/\mathfrak a)\geq 2$, then there exists two $\kk$-linearly independent element $c_1,\;c_2\in \mathrm{Soc}(P/\mathfrak a)$. Clearly, $ (\mathfrak a +(c_1))\cap (\mathfrak a +(c_2)) = \mathfrak a $, but this contradicts the fact that $\mathfrak a$ is irreducible, therefore $\dim_\kk \mathrm{Soc}(R/I)=1$ and $\mathfrak a$ is Gorenstein. 
\end{proof}

\begin{theorem}\label{maintheorem1}
    Let $(P,\m,\kk)$ be a regular local ring of dimension 3 and $I,\; J$ be two proper ideals in $P$. Then the product ideal $IJ$ is Golod.
\end{theorem}

\begin{proof}
     Set $R=P/IJ$, and $\textbf{K}^R=\mathbf{K}^R(\textbf{x})$. From [\cite{vandebogert2022products}, 3.1], we know that $H_1(\mathbf{K}^R)\cdot H_1(\mathbf{K}^R)=0$. It remains to prove $ H_1(\mathbf{K}^R)\cdot H_2(\mathbf{K}^R)=0$. 
     
     Suppose $H_1(\mathbf K^{R})\cdot H_2(\mathbf K^{R})\neq 0$. A basis for $ H_1(\mathbf K^{R})$ can be chosen to be represented by cycles of type $f\cdot v$, where $f\in I$ and $v\in \mathbf K^{R}_1$ is such that $\partial(v)\in J$. Therefore, we may choose $[z_1]=[fv]$ with $f\in I$ and $v\in \mathbf K^{R}_1$, and $[z_2]\in H_2(\mathbf K^{R})$ such that $[z_1\cdot z_2]\ne 0$. Similarly, a basis for $H_1(\mathbf K^{R})$ can be chosen to be represented by cycles of the form $gu$ with $g\in J$ and $u\in \mathbf K^{R}_1$ such that $\partial (u)\in I$. Therefore, we can also pick $[z_1']=[gu]$ with $g\in J$ and $u\in \mathbf K^{R}_1$ such that $[z_1'z_2]\ne 0$. 
     
     
     
     % pick $u,v\in \mathbf K^{R}_1$ and a cycle $z\in  \mathbf K^{R}_2$ such that $\partial(u)=f\in I$, $\partial(v)=g\in J$, and $[gu]\cdot [z]=[fv]\cdot [z]\neq 0$. 
     
     Since $K^{R}_3$ is a free module of rank one, we have  $[v\cdot z_2]=[rdx_{123}]$ and $[u\cdot z_2]=[sdx_{123}]$ for some $r, s\in R$.
     We claim that $r\notin J$ and $s\notin I$.
     Indeed, if $r\in J$, then $ [fv]\cdot [z_2]=[frdx_{123}]=0$, and if $s\in I$, then $ [gu]\cdot [z_2]=[gsdx_{123}]=0$, contradicting our assumption that $[z_1z_2]\ne 0$ and $[z_1'z_2]\ne 0$.
     
      From \cref{GorensteinEmbed} we have $\m$-primary Gorenstein ideals $\mathfrak a$ and $\mathfrak b$ such that $I\subseteq \mathfrak a$ with $s\notin \mathfrak a$ and $J\subseteq \mathfrak b$ with $r\notin \mathfrak b$. The following maps induced by the canonical projections are not zero:\[
    H_{2}(\mathbf{K}^R)\xrightarrow[]{\;\;\kappa_1\;\;}H_{2}(\mathbf{K}^{P/\mathfrak a}(\textbf{x}))\quad\text{and}\quad H_2(\mathbf{K}^R)\xrightarrow[]{\;\;\kappa_2\;\;}H_{2}(\mathbf{K}^{P/\mathfrak b}(\textbf{x})).
    \] 
    More precisely, we claim that  $ \kappa_1([z_2])\ne 0$ and $\kappa_2([z_2])\ne 0$. Using $\overline{\;w\;}$ to represent the image of an element $w\in\mathbf{K}^R$ along the canonical projection $\mathbf{K}^R\xrightarrow[]{} \mathbf{K}^{P/\mathfrak a}(\textbf{x})$, suppose by way of contradiction that $\overline{[z_2]}=0$. Then there exists $\ell\in \mathbf{K}^R_3$ such that $\partial(\overline{\ell})=\overline{z_2}$ in $ \mathbf{K}^{P/\mathfrak a}(\textbf{x})$. Since $\overline{u}\overline{\ell} \in \mathbf K_4^{P/\mathfrak a}=0$, we have $\partial(\overline{u})\overline{\ell}-\overline{u}\partial(\overline{\ell})=\partial(\overline{u}\overline{\ell})=0$. Recalling that $\partial(u)\in I\subseteq \mathfrak a$, which implies $\partial(\overline{u})=0$, we now have 
    $$
\overline{u}\overline{z_2}=\overline{u}\partial(\overline{\ell}) =\partial(\overline {u})\overline{\ell}=0.
    $$
    On the other hand, $[\overline{u}\overline{z_2}]=[\overline{sdx_{123}}],$ which is not zero since $s\notin \mathfrak a$. This is a contradiction.
    
    % Due to degree reasons, we have $$\partial(\overline{u}\overline{\ell})=0\implies \partial(\overline{u})\overline{\ell}-\overline{u}\partial(\overline{\ell})=0$$ 
    % Now, since $ \partial(\overline{u})=\overline{f}\in I\subseteq \mathfrak a$, and $\overline{u}\partial(\overline{\ell})=\overline{u}\overline{z_2}=\overline{sdx_{123}}$, we have $\partial(\overline{u})\overline{\ell}-\overline{u}\partial(\overline{\ell})=- \overline{sdx_{123}}=0$, but this is a contradiction as $ s\notin \mathfrak a$, therefore $\kappa_1$ is not a zero map, and similarly $\kappa_2$ is not a zero map.
    
    Now, \cref{mainlemma1} implies that there exists $\gamma \in \mathfrak a$ and $\beta \in \mathfrak b$ with $\mathfrak a\subseteq (\gamma)+\m J $ and $\mathfrak b\subseteq (\beta)+\m I $. This implies that $\mathfrak a+\mathfrak b = (\gamma,\beta)+\m( \mathfrak a+\mathfrak b)$. It follows by Nakayama's lemma that $\mathfrak a+\mathfrak b = (\gamma,\beta)$, which is a contradiction because $\mathfrak a + \mathfrak b$ is an $\m$-primary ideal and therefore it has grade 3.
\end{proof}

\begin{center}
    \section*{Acknowledgments}
\end{center}
    Experiments with Macaulay2 \cite{M2} helped us formulate results. 
    Some of the ideas in Section 6 grew out of a discussion with ChatGPT\cite{OpenAI2026ChatGPT}. We thank Liana \c{S}ega for helpful conversations.  
    
    % A preliminary version of Section 6, in which both ideals are Gorenstein, grew out of discussion with ChatGPT\cite{OpenAI2026ChatGPT}. Subsequently, working on those ideas, Author 1\dinesh{should I write author 1, or just we?} extended it to full generality. 


% \begin{theorem}\label{product of gorenstein}
%     Assume setting \ref{golodcriterionsetting} and additionally let $I$ be a grade 3 Gorenstein ideal.\begin{itemize}
%         \item[(a)] If $J\subseteq I$, then $IJ$ is Golod.
%         \item[(b)] If $J$ is also grade 3 Gorenstein, then $IJ$ is Golod. 
%     \end{itemize}
% \end{theorem}
% \begin{proof}
%     (a) Let $ J\subseteq I$, and suppose $\Tor_{2}^{P}(P/IJ,\kk)\xrightarrow[]{}\Tor_{2}^{P}(P/I,\kk)$ is non-zero map. From observation \ref{observation}, there exists $f\in I$ such that $I\subseteq (f)+\m J$. But $J\subseteq I$ implies that $ I=(f)+\m I$, and by graded Nakayama we have $I=(f)$, a contradiction as the grade of I is 3. Therefore, $\Tor_{2}^{P}(P/IJ,\kk)\xrightarrow[]{}\Tor_{2}^{P}(P/I,\kk)$ is zero map, and by \cref{GolodCriterion}, $IJ$ is Golod.   

%     (b) Let $J$ be a grade 3 Gorenstein. If both $\Tor_{2}^{P}(P/IJ,\kk)\xrightarrow[]{}\Tor_{2}^{P}(P/I,\kk)$ and $\Tor_{2}^{P}(P/IJ,\kk)\xrightarrow[]{}\Tor_{2}^{P}(P/J,\kk)$ is non-zero, then there exists $f\in I$ and $g\in J$, such that $I\subseteq (f)+\m J$ and $J\subseteq (g)+\m I$. So, we have $ I+J=(f,g)+\m(I+J)$. Now, graded Nakayama implies that $I+J= (f,g)$, a contradiction as the grade of $ I+J$ is 3, we have a contradiction. Therefore, one of the tor maps are zero and $IJ$ is Golod from \cref{GolodCriterion}.
% \end{proof}
% \begin{corollary}
%     Let $I$ be a grade 3 Gorenstein ideal in the regular local ring of embedding dimension 3. Then $I^k$ is Golod for all $k\geq 2$.
% \end{corollary}
% \begin{proof}
%     Let $k\geq 2$ and take $J=I^{k-1}$. The rest follows from part (a) of \cref{product of gorenstein}.
% \end{proof}


% \begin{theorem}
%     Assume setting \ref{golodcriterionsetting} and additionally let $I$ be a complete intersection. Then $IJ$ is Golod. 
% \end{theorem}

% \begin{proof}
%     Suppose $IJ$ is not Golod. Then $\Tor_{2}^{P}(P/IJ,\kk)\xrightarrow[]{}\Tor_{2}^{P}(P/I,\kk)$ is non-zero, and from observation \ref{observation}, one has $I\subseteq (f)+\m J$ where $I=(f,g,h)$. If $f\in \m J$, then $I \subseteq \m J$.
% \end{proof}

% From now on we will only be dealing with standard graded polynomial rings.
% \begin{definition}
% A standard graded Artinian $\kk$-algebra $R=\oplus_{i\geq 0}R_i$ with embedding dimension $e$, top socle degree s, and socle polynomial $\sum_{i=0}^s c_iz^i$ is called compressed if \[
% \dim_k R_i=\min\{\binom{e-1+i}{i},\sum_{j=i}^sc_P\binom{e-1+j-i}{j-i} \}
% \]
% for $i=0,\dots,s$.     
% \end{definition}
% In the standard graded polynomial ring $P=\kk[x_1,\cdots,x_n]$, we will call an ideal $I$ compressed if $R=P/I$ is Compressed. 
% \begin{remark}
%     \cref{GolodCriterion}, also hold when one replace regular local ring of embedding dimension 3, with a standard graded algebra of Krull dimension 3 and proper ideals with homogeneous proper ideals.
% \end{remark}
% \begin{proposition}\label{compressedproposition}
%     Let $P=\kk[x_1,x_2,x_3]$ be a standard graded algebra with homogeneous maximal ideal $\m$. Let $I\subseteq \m^2$ be a compressed ideal with initial degree a, top socle s, and socle polynomial $\sum_{i=0}^s c_iz^i$. Let $J$ be a proper homogeneous ideal with $J\subseteq \m^a$. Then $IJ$ is Golod.
% \end{proposition}
% \begin{proof}
%     Set $R=P/I$ and $\mathbf K^R=\mathbf{K}^{R}(\mathbf{x})$. Since $IJ\subseteq \m^{2a}$ we have for all $d\leq 2a$, $[\Tor_{2}^{P}(P/IJ,\kk)]_{d}= 0$. Since $P/I$ is compressed, we have $s\leq 2a-1$. Suppose $s< 2a-1$ and $d>2a$ or $s=2a-1$ and $d>2a+1$. In both situation we have $d-2>s$ and hence $[\Tor_{2}^{P}(R,\kk)]_{d}\cong [K_2^R]_d=\bigoplus_{1\leq i<j\leq 3}R_{d-2}dx_{ij}=0$. Suppose $s=2a-1$ and $d=2a+1$. This implies that $d=s+2$. Consider the map $\kk$-linear map \[
%     [K_3^R]_d\xrightarrow[]{\;\;\partial_3\;\;}[K_2^R]_d
%     \] 
%     Since $\ker(\partial_3)\cong(0:_R\m)_{s-1}$, we have $\dim \ker(\partial_3) =c_{s-1}$. Since $a\geq 2 $ we have $s-1\geq a$, which then implies that $\dim_{\kk} R_{d-3}=c_{s-1}+3c_P$. Now, by rank nullity, $\dim_{\kk} \mathrm{im}(\partial_3)=\dim_k [K_3^R]_d -\dim_k \ker(\partial_3)=\dim_{\kk} R_{d-3}-c_{s-1}=3c_s=\dim_{\kk} [K_2^R]_d$. Therefore $\partial_3$ is surjective, and $[\Tor_{2}^{P}(R,\kk)]_{d}=0$. Now, since Tor maps preserve internal degree we have $\Tor_{2}^{P}(P/IJ,\kk)\xrightarrow[]{}\Tor_{2}^{P}(P/I,\kk)$ is a zero map and from \cref{GolodCriterion}, $IJ$ is Golod.
% \end{proof}
% \cref{compressedproposition} yields a positive answer to the Puestion 4.11 proposed in \cite{vandebogert2022products}.
% \begin{corollary}
%     Let $P=\kk[x_1,x_2,x_3]$ be a standard graded ring with homogeneous maximal ideal $\m$. If $I$ and  $J$ be two compressed ideals in $P$, then their product $IJ$ is Golod. 
% \end{corollary}
% \begin{proof}
%     Suppose, $I,\;J\subseteq \m^2$. Let $a$ and $b$ be the initial degrees of $I$ and $J$, respectively. If $a\leq b$, then $J\subseteq \m^a$ and \cref{compressedproposition} implies $IJ$ is Golod. The case $a\geq b$ is the same. Now, suppose $I$ or $J$ contains a linear form; then the projective dimension of one of them is at most 2, and from \cref{Grade2Gorenstein}{} we have that $IJ$ is Golod. 
% \end{proof}

% \begin{corollary}
%     Let $I$ be a compressed ideal in $\kk[x_1,x_2,x_3]$. Then for $k\geq 2$, $I^k$ is Golod.   
% \end{corollary}
% \begin{proof}
%     If $I$ contains a linear form, then it follows from \cref{Grade2Gorenstein}, and if $I$ does not contain a linear form, then it follows from \cref{compressedproposition}.
    
% \end{proof}
%%%%%%%%%%%%%%%%%%%%%%%%%%%%%%%%%%%%%%%%%%%%%%%%%%%%%%%%%%%%%
%%%%%%%%%%%%%%%%%%%%%%%%%%%%%%%%%%%%%%%%%%%%%%%%%%%%%%%%%%%%


\bibliographystyle{abbrv}
\bibliography{references}

\end{document}